\documentclass[11pt]{article}
\usepackage[T1]{fontenc}
\usepackage{lmodern}
\usepackage[margin=1in]{geometry}
\usepackage{amsmath,amssymb,amsthm,mathtools}
\usepackage{booktabs,longtable,array}
\usepackage{graphicx,microtype,xcolor}
\usepackage{enumitem}
\usepackage{fancyhdr}
\usepackage[colorlinks=true,linkcolor=blue!45!black,citecolor=blue!45!black,urlcolor=blue!55!black]{hyperref}
\usepackage{xurl}
\definecolor{ink}{HTML}{16324F}
\setlength{\emergencystretch}{3em}
\setlength{\parskip}{0.2em}
\setlist{itemsep=0.25em,topsep=0.5em}
\pagestyle{fancy}
\fancyhf{}
\fancyhead[L]{\small\color{ink}Polylogarithms: exact structure and corrections}
\fancyhead[R]{\small ProveIt research continuation}
\fancyfoot[C]{\thepage}
\setlength{\headheight}{14pt}
\newtheorem{theorem}{Theorem}[section]
\newtheorem{proposition}[theorem]{Proposition}
\newtheorem{lemma}[theorem]{Lemma}
\newtheorem{corollary}[theorem]{Corollary}
\theoremstyle{definition}
\newtheorem{definition}[theorem]{Definition}
\theoremstyle{remark}
\newtheorem{remark}[theorem]{Remark}
\newcommand{\Li}{\operatorname{Li}}
\newcommand{\Cl}{\operatorname{Cl}}
\newcommand{\Log}{\operatorname{Log}}
\newcommand{\ii}{\mathrm{i}}
\newcommand{\Q}{\mathbb{Q}}
\newcommand{\C}{\mathbb{C}}
\newcommand{\Z}{\mathbb{Z}}
\newcommand{\N}{\mathbb{N}}
\newcommand{\Repart}{\operatorname{Re}}
\newcommand{\Impart}{\operatorname{Im}}
\newcommand{\Res}{\operatorname{Res}}
\newcommand{\sgn}{\operatorname{sgn}}
\newcommand{\fall}[2]{#1^{\underline{#2}}}
\newcommand{\src}[1]{\texttt{\small #1}}
\newcommand{\stirling}[2]{\genfrac{[}{]}{0pt}{}{#1}{#2}}
\DeclareMathOperator{\rank}{rank}
\DeclareMathOperator{\lcm}{lcm}
\DeclareMathOperator{\ord}{ord}
\hypersetup{pdftitle={Exact Structure and Certified Computation in the Polylogarithm Drafts},
pdfauthor={Research continuation prepared for Vladimir Reshetnikov with OpenAI assistance}}
\title{\color{ink}\textbf{Exact Structure and Certified Computation\\
in the Polylogarithm Drafts}\\[0.6em]
\large Mixed-point reductions, weighted distributions,\\
and Stieltjes derivative asymptotics}
\author{Research continuation for the ProveIt project\\
\small Prepared for Vladimir Reshetnikov with OpenAI assistance}
\date{October 8, 2026 (UTC)}
\begin{document}
\maketitle
\begin{abstract}
We continue the research in eight articles and associated reports in
ProveIt's polylogarithm documentation, using a fixed source snapshot.
The first result corrects a consequential structural claim: double
polylogarithms at inverse arguments reduce at every weight to the
already present one-direction doubles and single polylogarithms.
Four constants advertised as new mixed-point depth-two directions are
evaluated explicitly in depth-one terms. We also give a complete
even-weight formula and meromorphic generating function for the
alternating harmonic sums studied in the drafts.

For rational Hurwitz-zeta grids, an explicit character normal form proves
the all-denominator ranks suggested by the draft computations.
It remains valid over arbitrary commuting prime weights and finite jet
rings; resonance changes the adequacy of top-denominator coordinates,
not the rank of the full distribution module. For high parameter
derivatives of generalized Stieltjes constants, we prove a normally
convergent spectral expansion, exact rational error certificates at
parameter one, and a uniform reciprocal-gamma profile when the
Stieltjes index grows logarithmically with the derivative order.
Finally, we close the weight-two modular row-sum question with a
summation-order proof, all derivative rows, and explicit exponential
tail bounds.

The paper distinguishes classical ingredients and their specializations
from extensions proved here relative to the supplied corpus. It makes
no claim to prove period independence or a new transcendence theorem.
A source-specific correction register, reproducible computations, and
a prioritized research agenda accompany the proofs.
\end{abstract}
\vspace{0.5em}
\noindent\textbf{Source snapshot.}
\href{https://github.com/VladimirReshetnikov/ProveIt/commit/3a6d80ed618d839915deb0d19ab687f45e81d995}
{\texttt{3a6d80ed618d839915deb0d19ab687f45e81d995}}.
All source locators below refer to this commit. The original drafts
are not silently rewritten.
\clearpage
\tableofcontents
\clearpage
\section{Scope, contributions, and standards of proof}\label{sec:scope}

The source corpus connects cyclotomic multiple polylogarithms, rational
gamma and polygamma values, Hurwitz-zeta derivatives, generalized Stieltjes
constants, and modular lattice sums. Its strongest feature is the repeated
use of exact transformations across these subjects. Its main vulnerability
is the occasional promotion of a numerical search outcome to an
arithmetic or dimensional theorem.

An exact relation matrix determines the dimension of a \emph{formal
quotient by the rows actually imposed}. It does not establish that the
evaluated constants have that dimension. A failed integer-relation search
does not establish independence, even when repeated at higher precision.
Conversely, an explicit identity can disprove a proposed independence
claim without settling a transcendence conjecture.

The present continuation has five strands.
\begin{enumerate}
\item Section~\ref{sec:mixed} reduces
\(\Li_{a,b}(\xi,\xi^{-1})\) at every weight to the existing
\((\xi,1)\) family and evaluates the four named mixed-point survivors.
\item Section~\ref{sec:euler} proves the complete even-weight evaluation
of \(S_p=\sum_{n\ge0}(-1)^nH_n/(2n+1)^p\), and determines its even
meromorphic generating function, whose first poles occur at \(\pm3\).
\item Sections~\ref{sec:distribution}--\ref{sec:bridges} prove the
all-denominator distribution ranks and give a single local identity for
all positive/negative \(L\)-function derivative bridges.
\item Section~\ref{sec:asymptotics} provides quantitative control of high
parameter derivatives of Stieltjes constants: a spectral expansion,
two uniform asymptotic terms, and exact rational enclosures.
\item Section~\ref{sec:modular} treats weight-two modular row sums,
where summation order matters, and gives derivative and tail formulas.
\end{enumerate}

\subsection{Relation to existing literature}

These are extensions and repairs of the supplied research programme.
The mixed parity reductions are special cases of Panzer's general parity
theorem~\cite{Panzer}. The odd-index \(S_p\) evaluations specialize the
Euler-sum theory of Xu and Wang~\cite{XuWang}; we supply an independent
beta-integral proof in the repository's normalization. The rank principle
belongs to the classical theory of universal
distributions~\cite{Kubert,Ouyang}; our elementary normal form makes the
Stieltjes applications and coordinate resonances explicit. The original
all-order Stieltjes parameter-derivative identity is due to
Coffey~\cite{Coffey}.

The reciprocal-gamma factor in Section~\ref{sec:asymptotics} is related to
classical first-kind Stirling-number and permutation-cycle asymptotics;
see Kowalski and Nikeghbali~\cite{KowalskiNikeghbali}.
The contribution here is its quantitative transfer to the parameter
derivative tower, including a uniform second term and rational
certificates. The weight-two anomaly and the differential algebra of
Eisenstein series are classical~\cite{RomikScherer,Ramanujan}.
We do not claim priority for these general mechanisms. We give complete
proofs of the specific extensions and state precisely what they settle.

\subsection{Notation and source coverage}

Our convention is
\[
\Li_{a,b}(u,v)=\sum_{n>m\ge1}\frac{u^n v^m}{n^a m^b}.
\]
When comparing with an increasing-index convention, reverse both
the weight list and the argument list. We distinguish
\[
\omega=e^{2\pi\ii/3},\qquad \tau_{\mathrm{hex}}=e^{\pi\ii/3}.
\]
The first is a cube root; the second is the hexagonal lattice parameter.
We write
\[
\beta(s)=\sum_{n\ge0}\frac{(-1)^n}{(2n+1)^s},\qquad
C_{2j}=\Impart\Li_{2j}(\omega)=\Cl_{2j}(2\pi/3),\qquad G=\beta(2).
\]
No irrationality assumption about \(G\) or the other even beta values
is made.

All eight article sources and the report inventory were inspected.
Detailed proof auditing concentrates on the multiple-polylog, Stieltjes,
gamma-relation, and modular-row strands used here.
We do not claim an exhaustive independent audit of every formula in the
Herglotz and other reports. Appendix~\ref{app:audit} identifies the
affected sources and distinguishes existing README warnings from
corrections established in this continuation.

\section{Inverse-argument mixed doubles add no new directions}
\label{sec:mixed}

Put \(M_{a,b}(z)=\Li_{a,b}(z,z^{-1})\) and \(w=a+b\).
We allow \(0<|z|<1\), and \(|z|=1,z\ne1\), with positive integers
\(a,b\). Boundary sums use the stated outer-index order. They converge
absolutely if \(a\ge2\). For \(a=1\), the inner colored harmonic sum
equals its limit plus \(O(n^{-b})\); the outer sum is therefore a
convergent oscillatory series plus an absolutely convergent remainder.

\begin{theorem}[All-weight inverse-argument reduction]\label{thm:mixed}
In this domain,
\begin{align}
M_{a,b}(z)
={}&(-1)^{b+1}\sum_{j=1}^{a}
 \binom{w-j-1}{b-1}\Li_{w-j,j}(z,1)
 +(-1)^{b+1}\binom{w-1}{b}\Li_w(z)\notag\\
&+\sum_{j=2}^{b}(-1)^{b-j}
 \binom{w-j-1}{a-1}\zeta(j)\Li_{w-j}(z)\notag\\
&+(-1)^b\sum_{j=2}^{a}
 \binom{w-j-1}{b-1}\zeta(j)\Li_{w-j}(z).
\label{eq:mixedgeneral}
\end{align}
Empty sums are zero. In particular, \(F_a=M_{a,1}\) satisfies
\begin{equation}
F_a(z)=a\Li_{a+1}(z)+\sum_{j=1}^{a}\Li_{a+1-j,j}(z,1)
       -\sum_{j=2}^{a}\zeta(j)\Li_{a+1-j}(z).
\label{eq:mixedb1}
\end{equation}
\end{theorem}

\begin{proof}
Set \(h=n-m\). Principal parts at \(m=0,-h\) give
\begin{align}
\frac1{m^b(m+h)^a}
={}&\sum_{j=1}^{b}
 \frac{(-1)^{b-j}\binom{w-j-1}{a-1}}{h^{w-j}m^j}\notag\\
&+(-1)^b\sum_{j=1}^{a}
 \frac{\binom{w-j-1}{b-1}}{h^{w-j}(m+h)^j}.
\label{eq:partialfraction}
\end{align}
Indeed, the difference has no poles and tends to zero at infinity.
The \(j=1\) coefficients are opposite; their sum over \(m\) is finite
because \(\sum_{m\ge1}(1/m-1/(m+h))=H_h\).
For \(j\ge2\), the shifted sum is \(\zeta(j)-H_h^{(j)}\).
Consequently
\begin{align*}
c_h:=\sum_{m\ge1}\frac1{m^b(m+h)^a}
={}&(-1)^{b+1}\sum_{j=1}^{a}
 \binom{w-j-1}{b-1}\frac{H_h^{(j)}}{h^{w-j}}\\
&+\sum_{j=2}^{b}(-1)^{b-j}\binom{w-j-1}{a-1}
   \frac{\zeta(j)}{h^{w-j}}
 +(-1)^b\sum_{j=2}^{a}\binom{w-j-1}{b-1}
   \frac{\zeta(j)}{h^{w-j}}.
\end{align*}
Multiply by \(z^h\), sum, and use
\[
\sum_{h\ge1}\frac{z^h H_h^{(j)}}{h^{w-j}}
=\Li_{w-j,j}(z,1)+\Li_w(z),\qquad
\sum_{j=1}^{a}\binom{w-j-1}{b-1}=\binom{w-1}{b}.
\]
For \(|z|<1\) these operations are absolute.

To justify the boundary even when \(a=1\), let \(O_N\) be the original
sum over \(n\le N\) and \(D_N=\sum_{h=1}^{N-1}z^h c_h\).
Partial summation gives
\[
\left|\sum_{h=h_0}^{h_1}\frac{z^h}{(m+h)^a}\right|
\le \frac{4}{|1-z|(m+h_0)^a}.
\]
Splitting the discrepancy at \(m=N\) now yields
\[
|D_N-O_N|
\le\frac4{|1-z|}
\left(N^{-a}H_N^{(b)}+\sum_{m>N}m^{-a-b}\right)\longrightarrow0.
\]
The coefficients \(c_h\) decrease to zero, the latter by dominated
convergence with majorant \(m^{-a-b}\).
Dirichlet's test proves convergence of \(\sum z^h c_h\).
The one-direction doubles on the right also converge. This proves the
identity in the original boundary convention.
\end{proof}

\begin{corollary}[The mixed quotient is zero]\label{cor:mixedzero}
Let \(\mathcal A_w(z)\) be the \(\Q\)-span of
\[
\Li_{r,w-r}(z,1)\ (1\le r<w),\quad
\Li_w(z),\quad \zeta(j)\Li_{w-j}(z)\ (2\le j<w).
\]
Every \(M_{a,b}(z)\) of weight \(w\) belongs to \(\mathcal A_w(z)\).
Adding the mixed values contributes zero dimension modulo this space.
\end{corollary}

At \(z=\ii\), this is the Gaussian family already used in the report.
At \(z=\omega^2\), the one-direction values are conjugates of its
\((\omega,1)\) values, whose real and imaginary parts it includes.
Thus the proposed additional mixed directions disappear at every
weight. The result concerns inverse-color pairs, and does not assert
that all colored double polylogarithms reduce to one direction.

\subsection{An independent integral representation}

\begin{proposition}\label{prop:mixedintegral}
In the same domain,
\begin{align}
M_{a,b}(z)&=\frac{z}{(a-1)!}\int_0^1
 \frac{(-\log t)^{a-1}\Li_b(t)}{1-zt}\,dt,
\label{eq:mixedintegral}\\
\Li_{a,b}(z,1)&=\frac{z}{(a-1)!}\int_0^1
 \frac{(-\log t)^{a-1}\Li_b(zt)}{1-zt}\,dt.
\label{eq:directionintegral}
\end{align}
\end{proposition}
\begin{proof}
Insert
\(n^{-a}=(a-1)!^{-1}\int_0^1(-\log t)^{a-1}t^{n-1}\,dt\).
The \(h=n-m\) geometric sum gives \(zt/(1-zt)\);
the \(m\)-sum gives \(\Li_b(t)\) or \(\Li_b(zt)\).
At a fixed boundary point \(z\ne1\), the denominator is bounded away
from zero on \([0,1]\), and the possible logarithmic endpoint
singularities are integrable. Dominated convergence and the preceding
boundary argument complete the proof.
\end{proof}

These formulas supply independent numerical checks without the
power-law acceleration problem at \(z_1z_2=1\). Quadrature itself is
not a rigorous certificate unless its numerical error is enclosed.

\subsection{The mixed parity component reduces to single values}

Panzer's theorem predicts lower depth for the real component at odd
weight and the imaginary component at even weight~\cite{Panzer}.
Here is a direct proof for the mixed \(F_a\) family.

\begin{theorem}[Mixed parity formulas]\label{thm:mixedparity}
For \(|z|=1,z\ne1\) and \(m\ge1\),
\begin{align}
\Repart F_{2m}(z)
={}&\frac{2m+1}{2}\Repart\Li_{2m+1}(z)
 -\sum_{j=1}^{m}\zeta(2j)\Repart\Li_{2m+1-2j}(z)\notag\\
&-\frac12\zeta(2m+1)
 +\Repart\Li_{2m}(z)\Repart\Li_1(z).
\label{eq:realparity}
\end{align}
For \(m\ge0\),
\begin{align}
\Impart F_{2m+1}(z)
={}&(m+1)\Impart\Li_{2m+2}(z)
 -\sum_{j=1}^{m}\zeta(2j)\Impart\Li_{2m+2-2j}(z)\notag\\
&+\Impart\Li_{2m+1}(z)\Repart\Li_1(z).
\label{eq:imagparity}
\end{align}
\end{theorem}
\begin{proof}
For \(a\ge2\), set
\(\mathcal B_a(z)=-\Li_a(z)-(-1)^a\Li_a(z^{-1})\).
The Fourier series
\(-\mathcal B_a(z)=\sum_{n\ne0}z^n/n^a\) is absolutely summable,
while \(\Li_1(z^{-1})=\sum_{\ell\ge1}z^{-\ell}/\ell\) belongs to
\(L^2\) on the circle. Their product has positive \(h\)-th coefficient
\(c_h=\sum_{\ell\ge1}[\ell(\ell+h)^a]^{-1}\), constant coefficient
\(\zeta(a+1)\), and negative \(h\)-th coefficient
\[
q_{-h}=\sum_{\ell\ge1,\ \ell\ne h}\frac1{\ell(\ell-h)^a}.
\]
The difference \(q_{-h}-(-1)^a c_h\) is
\begin{equation}
\sum_{\ell\in\Z\setminus\{0,h\}}\frac1{\ell(\ell-h)^a}
=(-1)^a\left[
 2\sum_{j=1}^{\lfloor a/2\rfloor}
       \frac{\zeta(2j)}{h^{a+1-2j}}
 -\frac{a+1}{h^{a+1}}\right].
\label{eq:bilateralpartial}
\end{equation}
To check this absolutely convergent bilateral sum, use
\[
\frac1{\ell(\ell-h)^a}
=\frac1{(-h)^a\ell}
 -\sum_{r=1}^{a}\frac1{(-h)^{a+1-r}(\ell-h)^r}.
\]
Symmetric summation of the \(r=1\) terms gives respectively
\(-1/h\) and \(1/h\). For \(r\ge2\), the shifted sum is
\((1+(-1)^r)\zeta(r)-(-h)^{-r}\). Substitution gives
\eqref{eq:bilateralpartial}.
Comparing all Fourier coefficients now gives
\begin{align}
F_a(z)+(-1)^aF_a(z^{-1})
={}&(-1)^a\left[(a+1)\Li_{a+1}(z^{-1})
 -2\sum_{j=1}^{\lfloor a/2\rfloor}
       \zeta(2j)\Li_{a+1-2j}(z^{-1})\right]\notag\\
&-\zeta(a+1)-\Li_1(z^{-1})\mathcal B_a(z).
\label{eq:mixedfourier}
\end{align}
The equality holds first in \(L^2\), then pointwise by continuity away
from \(z=1\). Real parts for even \(a\), and imaginary parts for odd
\(a\), give the claims. The remaining case \(a=1\) follows from
\(F_1(z)=\Li_2(z)+\tfrac12\Li_1(z)^2\), itself a consequence of
\eqref{eq:mixedb1} and the elementary shuffle identity.
\end{proof}

\subsection{Four explicit replacements for the claimed survivors}

\begin{corollary}\label{cor:four}
The four highlighted mixed-point constants satisfy
\begin{align}
\Repart\Li_{4,1}(\ii,-\ii)
&=-\frac{587}{1024}\zeta(5)
 +\frac3{32}\zeta(2)\zeta(3)
 +\frac{135}{256}\zeta(4)\log2,
\label{eq:gaussian4}\\
\Impart\Li_{5,1}(\ii,-\ii)
&=3\beta(6)-\zeta(2)\beta(4)-\zeta(4)G
 -\frac{5\pi^5\log2}{3072},
\label{eq:gaussian5}\\
\Repart\Li_{4,1}(\omega^2,\omega)
&=-\frac{281}{162}\zeta(5)
 +\frac49\zeta(2)\zeta(3)
 +\frac{20}{27}\zeta(4)\log3,
\label{eq:eisenstein4}\\
\Impart\Li_{5,1}(\omega^2,\omega)
&=-3C_6+\zeta(2)C_4+\zeta(4)C_2
 +\frac{\pi^5\log3}{729}.
\label{eq:eisenstein5}
\end{align}
\end{corollary}
\begin{proof}
Use Theorem~\ref{thm:mixedparity} with \(m=2\), together with
\[
\Repart\Li_s(\ii)=-2^{-s}(1-2^{1-s})\zeta(s),\qquad
\Repart\Li_s(\omega)=\frac{3^{1-s}-1}{2}\zeta(s)\quad(s>1),
\]
and
\[
\Repart\Li_1(\ii)=-\tfrac12\log2,\quad
\Repart\Li_1(\omega)=-\tfrac12\log3,\quad
\Impart\Li_5(\ii)=\frac{5\pi^5}{1536},\quad
\Impart\Li_5(\omega)=\frac{2\pi^5}{729}.
\]
The real identities follow by residue-class splitting; the imaginary
fifth-order values follow from the odd-beta and Bernoulli Fourier
formulas, with \(B_5(1/3)=-5/243\).
Conjugation changes the sign of imaginary parts at \(\omega^2\).
\end{proof}

The four values, in the displayed order, are approximately
\[
\begin{aligned}
&-0.01342026735728779,\qquad && 0.03269424445221378,\\
&-0.03903921845650667,\qquad &&-0.02533938136807792.
\end{aligned}
\]
Their right sides already belong to the report's listed single-value
bases. Its \src{thm:parallel} is therefore false in the printed
convention. The evaluator scripts and value stores cited by that report
were not imported into ProveIt, so their precise failure cannot be
diagnosed from the available files. Replace the survivor claims by
Theorem~\ref{thm:mixed} and Corollary~\ref{cor:four}, and recompute any
dimension conclusions that depended on the proposed new directions.

\section{Alternating harmonic sums: complete even-weight evaluation}
\label{sec:euler}

The Gaussian depth draft studies
\begin{equation}
S_p=\sum_{n=0}^{\infty}\frac{(-1)^nH_n}{(2n+1)^p},
\qquad H_0=0,\qquad H_n=\sum_{k=1}^{n}\frac1k.
\label{eq:euler-definition}
\end{equation}
For ordinary convergence the correct index range is \(p\ge1\).
At \(p=0\), the summands do not tend to zero. An Abel value at that
index would be a separate definition, not the sum in
\eqref{eq:euler-definition}.

The odd-index evaluations below are a specialization of the existing
Euler-sum literature. We prove them independently because the integral
argument explains exactly which symmetry supplies the evaluation and
leads to a useful global generating function. Neither the evaluation
nor the symmetry supplies a nonmembership theorem for the even indices.

\subsection{Convergence and the Mellin representation}

\begin{lemma}\label{lem:euler-mellin}
For every integer \(p\ge1\), the series \eqref{eq:euler-definition}
converges and satisfies \(|S_p|\le3^{-p}\). It converges absolutely
precisely when \(p>1\). With
\[
Q(x)=\frac{\log(1+x^2)}{1+x^2},
\]
one has
\begin{equation}
S_p=-\frac1{(p-1)!}\int_0^1(-\log x)^{p-1}Q(x)\,dx.
\label{eq:euler-mellin}
\end{equation}
\end{lemma}

\begin{proof}
For \(n\ge1\),
\[
\frac{H_{n+1}}{H_n}
=1+\frac1{(n+1)H_n}
\le1+\frac1{n+1}
<\frac{2n+3}{2n+1}.
\]
Thus \(H_n/(2n+1)^p\) decreases to zero for every \(p\ge1\).
The alternating-series theorem gives convergence and the bound by
the first nonzero term, \(3^{-p}\). The estimate
\(H_n\le1+\log n\) proves absolute convergence for \(p>1\);
for \(p=1\), even the lower bound \(H_n\ge1\) proves divergence
of the absolute series.

To justify the endpoint \(p=1\) together with the other indices, take
\(0\le r<1\). The harmonic-number generating identity gives
\[
\sum_{n\ge0}(-r)^nH_nx^{2n}
=-\frac{\log(1+rx^2)}{1+rx^2},\qquad 0\le x\le1.
\]
The series is uniformly absolutely convergent in \(x\), and
\[
\int_0^1x^{2n}(-\log x)^{p-1}\,dx
=\frac{(p-1)!}{(2n+1)^p}.
\]
Termwise integration therefore proves the regularized version of
\eqref{eq:euler-mellin}. Since
\[
0\le \frac{\log(1+rx^2)}{1+rx^2}\le x^2,
\]
dominated convergence gives the stated integral as \(r\uparrow1\).
Abel's theorem identifies the series limit with the already established
ordinary sum. In particular, no unproved interchange of a conditionally
convergent boundary series is required.
\end{proof}

\subsection{Evaluation at every odd index}

\begin{theorem}[All odd-index harmonic sums]\label{thm:euler-odd}
For every integer \(m\ge0\),
\begin{equation}
\boxed{\displaystyle
S_{2m+1}=(2m+1)\beta(2m+2)-2\beta(2m+1)\log2
-\sum_{j=1}^{m}(2-2^{-2j})
 \beta(2m-2j+1)\zeta(2j+1).}
\label{eq:euler-odd}
\end{equation}
The sum is empty for \(m=0\).
\end{theorem}

\begin{proof}
We first relate a half-line logarithmic moment to a full-line moment.
Put
\[
J_m=\int_0^1\log^{2m}x\,Q(x)\,dx,
\qquad
L_m=\int_0^{\infty}\log^{2m}x\,Q(x)\,dx.
\]
All these integrals converge: \(Q(x)=O(x^2)\) at zero and
\(Q(x)=O(x^{-2}\log x)\) at infinity.
Substitution \(x=1/u\) in the part over \((1,\infty)\) uses the
exact identity
\[
\frac{Q(1/u)}{u^2}
=Q(u)-\frac{2\log u}{1+u^2}.
\]
Consequently,
\begin{align}
L_m
&=2J_m-2\int_0^1\frac{\log^{2m+1}u}{1+u^2}\,du\notag\\
&=-2(2m)!S_{2m+1}+2(2m+1)!\beta(2m+2).
\label{eq:euler-inversion}
\end{align}
Here the first term follows from Lemma~\ref{lem:euler-mellin}; the
second follows by expanding \((1+u^2)^{-1}\). That integration is
absolutely summable, since the integrated terms have size
\((2m+1)!/(2n+1)^{2m+2}\).

It remains to evaluate \(L_m\). Define the full Mellin transform
\[
F(a)=\int_0^{\infty}x^{a-1}Q(x)\,dx.
\]
The integral in fact converges for \(-2<\Repart a<2\). On the
smaller strip \(0<\Repart a<2\), the beta integral gives
\begin{equation}
F(a)=\frac{\pi}{2\sin(\pi a/2)}
       \bigl[\psi(1)-\psi(1-a/2)\bigr],
\label{eq:euler-full-mellin}
\end{equation}
where \(\psi=\Gamma'/\Gamma\). Indeed, for
\(\Repart b>\Repart a/2>0\),
\[
\int_0^{\infty}x^{a-1}(1+x^2)^{-b}\,dx
=\frac{\Gamma(a/2)\Gamma(b-a/2)}{2\Gamma(b)}.
\]
Differentiate with respect to \(b\), negate the result, and set
\(b=1\). Differentiation under the integral is valid uniformly on
compact subsets of this parameter region, because the additional
factor \(\log(1+x^2)\) preserves the integrable bounds at zero
and infinity. Euler's reflection formula then yields
\eqref{eq:euler-full-mellin}. The same bounds with further powers
of \(\log x\) justify differentiation in \(a\), so
\begin{equation}
F^{(2m)}(1)=L_m.
\label{eq:euler-moments}
\end{equation}

The digamma duplication formula and the convergent polygamma series
give
\[
\psi(1)-\psi(1/2)=2\log2,
\qquad
\psi^{(k)}(1/2)
=(-1)^{k+1}k!(2^{k+1}-1)\zeta(k+1),\quad k\ge1.
\]
For clarity, the second identity follows by separating the odd terms
in \(\psi^{(k)}(z)=(-1)^{k+1}k!\sum_{n\ge0}(n+z)^{-k-1}\).
Expanding \eqref{eq:euler-full-mellin} at \(a=1\), its even part
is therefore
\begin{equation}
\frac{F(1+t)+F(1-t)}2
=\pi\sec\!\left(\frac{\pi t}{2}\right)
 \left[\log2+\sum_{j\ge1}
 (1-2^{-2j-1})\zeta(2j+1)t^{2j}\right],
\qquad |t|<1.
\label{eq:euler-even-mellin}
\end{equation}

The needed coefficients of the secant can be obtained without any
assumption about special values. Another beta integral and the same
inversion give
\[
\frac{\pi}{2}\sec\!\left(\frac{\pi t}{2}\right)
=\int_0^1\frac{x^t+x^{-t}}{1+x^2}\,dx
=2\sum_{r\ge0}\beta(2r+1)t^{2r},\qquad |t|<1.
\]
To justify the last equality, expand the numerator in powers of \(t\)
under an integrable bound \(x^{-|t|}\). The logarithmic moments
with \(r\ge1\) may be integrated termwise absolutely; for \(r=0\)
the identity is \(\arctan1=\beta(1)=\pi/4\).
Taking the coefficient of \(t^{2m}\) in
\eqref{eq:euler-even-mellin} now yields
\[
\frac{L_m}{(2m)!}
=4\beta(2m+1)\log2
 +4\sum_{j=1}^{m}(1-2^{-2j-1})
       \beta(2m-2j+1)\zeta(2j+1).
\]
Substituting this into \eqref{eq:euler-inversion} and dividing by
\(2(2m)!\) proves \eqref{eq:euler-odd}.
\end{proof}

\paragraph{Literature identification and normalization.}
In Xu and Wang's notation for linear Euler \(R\)-sums,
\[
R_{1,q}^{1,-1}
=\sum_{n\ge1}\frac{(-1)^{n-1}H_{n-1}}{(n-1/2)^q}
=2^qS_q.
\]
For odd \(q\ge3\), equation~(20) on printed page~994 of
\cite{XuWang} specializes to \eqref{eq:euler-odd}.
The conversion uses their conventions
\[
\widetilde t(k;-1)=-2^k\beta(k),\qquad
\widetilde t(k;1)=(2^k-1)\zeta(k)\quad(k\ge2),
\qquad \widetilde t(1;1)=2\log2.
\]
The \(m=0\) case follows directly from the proof above. Thus the
theorem is an explicit, independently proved specialization of known
results, not a claim of a globally new evaluation theorem.

For example,
\begin{align*}
S_1&=G-\frac\pi2\log2,\\
S_3&=3\beta(4)-\frac{\pi^3}{16}\log2
                      -\frac{7\pi}{16}\zeta(3),\\
S_9&=9\beta(10)-\frac{1385\pi^9}{20643840}\log2
       -\frac{427\pi^7}{737280}\zeta(3)
       -\frac{155\pi^5}{24576}\zeta(5)\\
   &\hspace{1.7em}-\frac{127\pi^3}{2048}\zeta(7)
       -\frac{511\pi}{1024}\zeta(9).
\end{align*}
Every odd beta value is a rational multiple of the corresponding odd
power of \(\pi\), as also follows by comparing the secant coefficients
above. Hence \eqref{eq:euler-odd} is homogeneous of weight \(2m+2\)
in the single-value generators used in the draft. In particular, it
proves the positive, even-weight half of the proposed alternation
pattern at all weights.

\subsection{A meromorphic generating function and its exact singularities}

\begin{theorem}[Even generator]\label{thm:euler-generator}
The power series
\[
E(t)=\sum_{m\ge0}S_{2m+1}t^{2m}
\]
has radius of convergence exactly \(3\). It extends to an even
meromorphic function on \(\C\), with the closed form
\begin{equation}
\boxed{\displaystyle
E(t)=\frac{\psi_1((1+t)/4)-\psi_1((3+t)/4)}{16}
 +\frac{\pi\,[\psi((1+t)/2)-\psi(1)]}
        {4\cos(\pi t/2)},}
\label{eq:euler-generator}
\end{equation}
where \(\psi_1=\psi'\), and apparent singularities are interpreted
by meromorphic continuation. Its only poles are simple and occur at
\(t=\pm(2n+1)\), \(n\ge1\), with
\begin{equation}
\Res_{t=2n+1}E(t)=\frac{(-1)^{n+1}H_n}{2},\qquad
\Res_{t=-(2n+1)}E(t)=\frac{(-1)^nH_n}{2}.
\label{eq:euler-residues}
\end{equation}
In particular, the apparent singularities at \(t=\pm1\) are
removable, and
\begin{equation}
E(1)=E(-1)=-\frac{\pi^2}{48}.
\label{eq:euler-removable}
\end{equation}
\end{theorem}

\begin{proof}
Lemma~\ref{lem:euler-mellin} and the bound \(|S_p|\le3^{-p}\)
first define \(E\) on \(|t|<3\). Summing the even exponential
series under the Mellin integral gives
\begin{equation}
E(t)=-\frac12\int_0^1(x^t+x^{-t})Q(x)\,dx.
\label{eq:euler-generator-integral}
\end{equation}
For \(|t|<3\), absolute values of the summed integrand are bounded
by \(Q(x)\cosh(|t|\log x)\), which is integrable at zero since
\(Q(x)=O(x^2)\). The right-hand side of
\eqref{eq:euler-generator-integral} is holomorphic on the larger
vertical strip \(|\Repart t|<3\): on each compact subset, the same
argument applies with a fixed exponent strictly smaller than \(3\),
and all parameter derivatives add only powers of \(|\log x|\).
It also proves evenness there.

To evaluate the integral, let
\[
J(t)=\int_0^1x^tQ(x)\,dx,
\qquad
K(t)=\int_0^1\frac{x^t\log x}{1+x^2}\,dx.
\]
For \(|\Repart t|<1\), splitting the full Mellin integral at one
and inverting the second part gives
\[
F(1-t)=J(-t)+J(t)-2K(t).
\]
In this strip, expansion of \((1+x^2)^{-1}\) and absolute
summability after integration give
\begin{align*}
K(t)
&=-\sum_{n\ge0}\frac{(-1)^n}{(2n+1+t)^2}\\
&=-\frac1{16}
 \bigl[\psi_1((1+t)/4)-\psi_1((3+t)/4)\bigr].
\end{align*}
Since \(E(t)=-F(1-t)/2-K(t)\), substitution of
\eqref{eq:euler-full-mellin} proves
\eqref{eq:euler-generator} in \(|\Repart t|<1\).

The exact global singularity statement should not be inferred merely
from the individual terms of that formula, because they have cancelling
poles. Instead, retain the conditionally convergent coefficient
\(S_1\) separately and sum all the absolutely convergent higher
coefficients. For \(|t|<3\), this gives
\begin{equation}
E(t)=S_1+\sum_{n\ge1}
 \frac{(-1)^nH_n t^2}
 {(2n+1)((2n+1)^2-t^2)}.
\label{eq:euler-generator-partial-fractions}
\end{equation}
The interchange is valid because
\[
\sum_{n\ge1}\frac{H_n}{2n+1}
 \sum_{m\ge1}\left(\frac{|t|}{2n+1}\right)^{2m}
=\sum_{n\ge1}
 \frac{H_n|t|^2}{(2n+1)((2n+1)^2-|t|^2)}<\infty.
\]
For any compact set avoiding the displayed poles in
\eqref{eq:euler-generator-partial-fractions}, the summands for
sufficiently large \(n\) are uniformly
\(O((1+\log n)/n^3)\). Thus the series converges normally there
and defines a meromorphic function on \(\C\). At a point
\(2n+1\) or \(-(2n+1)\), precisely one summand is singular;
its elementary residue is \eqref{eq:euler-residues}. These
residues are nonzero, and there are no further poles. The right-hand
side of \eqref{eq:euler-generator} is also meromorphic on \(\C\)
and agrees on an open strip, so the meromorphic identity theorem
proves that the two representations agree everywhere. In particular,
all additional apparent poles in the polygamma expression cancel.
The nearest actual poles to zero are \(3\) and \(-3\), proving
that the Taylor radius is exactly \(3\).

Finally, \eqref{eq:euler-generator-integral} can be evaluated directly
at \(t=1\), inside its holomorphy strip:
\[
E(1)=-\frac12\int_0^1\frac{\log(1+x^2)}{x}\,dx
     =\frac14\Li_2(-1)=-\frac{\pi^2}{48}.
\]
Evenness gives the same value at \(-1\), proving
\eqref{eq:euler-removable} without a delicate cancellation of
separate Laurent expansions.
\end{proof}

\subsection{The full generator and the unresolved parity direction}

The symmetry can be stated more precisely by keeping all indices.

\begin{proposition}\label{prop:euler-full-generator}
The full generating function
\[
A(t)=\sum_{p\ge1}S_pt^{p-1}
\]
has Taylor radius \(3\), integral representation
\begin{equation}
A(t)=-\int_0^1x^{-t}Q(x)\,dx,\qquad \Repart t<3,
\label{eq:euler-full-generator-integral}
\end{equation}
and global meromorphic representation
\begin{equation}
A(t)=S_1+\sum_{n\ge1}
 \frac{(-1)^nH_n t}{(2n+1)(2n+1-t)}.
\label{eq:euler-full-generator-poles}
\end{equation}
Its only poles are simple, at \(t=2n+1\), \(n\ge1\), with
residue \((-1)^{n+1}H_n\). Moreover,
\begin{equation}
\frac{A(t)+A(-t)}2=E(t),\qquad
\frac{A(t)-A(-t)}2=\sum_{m\ge1}S_{2m}t^{2m-1}
\quad (|t|<3).
\label{eq:euler-parity-generators}
\end{equation}
\end{proposition}

\begin{proof}
For \(|t|<3\), insert \eqref{eq:euler-mellin} and sum the
exponential series \(\sum_{p\ge1}t^{p-1}(-\log x)^{p-1}/(p-1)!
=x^{-t}\). The absolute sum is bounded by \(x^{-|t|}Q(x)\),
which is integrable. The resulting integral is holomorphic for
\(\Repart t<3\), by the same compact-set bounds as above.
Separating \(p=1\) and summing the geometric series for \(p\ge2\)
gives \eqref{eq:euler-full-generator-poles}; the needed absolute
bound is
\[
\sum_{n\ge1}\frac{H_n|t|}{(2n+1)(2n+1-|t|)}<\infty,
\qquad |t|<3.
\]
On a compact set avoiding positive odd integers at least three,
the resulting rational summands are uniformly
\(O((1+\log n)/n^2)\). Normal convergence gives the stated global
continuation. The single summand singular at \(t=2n+1\) has residue
\((-1)^{n+1}H_n\), and that residue is nonzero. This proves both
the complete pole list and the exact radius. Taking even and odd
parts of the Taylor series proves \eqref{eq:euler-parity-generators}.
\end{proof}

The inversion \(x\mapsto1/x\) determines the even part of \(A\)
through the full Mellin transform. It leaves the odd part as a
different analytic quantity. The absence of the same reduction for
that part is not a proof that its coefficients lie outside the algebra
generated by single polylogarithmic values. In particular, no result
here proves nonmembership for \(S_{2m}\), which have odd total
weight. Such a claim requires an additional arithmetic or motivic
argument for the specified constants and specified algebra. The
presence of an even beta value likewise supplies no independence
certificate: even the irrationality of \(G=\beta(2)\) is not known.
The correction to the draft is therefore precise: replace the proposed
two-sided alternation theorem by \eqref{eq:euler-odd} and its proved
even-weight membership statement, and retain the opposite direction
as a research question.

\section{Weighted distribution modules and exact rational-grid ranks}
\label{sec:distribution}

The distribution relations in the Stieltjes drafts have a simple
representation-theoretic structure. The relevant dimension is the
dimension of a formal quotient by specified relations. The classical
theory of universal distributions provides the broader setting
\cite{Kubert,Ouyang}. Here we give an explicit character normal form,
including its behavior at exceptional weights and in finite Taylor
jets. This proves the two all-denominator rank statements suggested by
the source computations, without making an independence assertion about
evaluated special values.

\subsection{The relation module and its complete character decomposition}

Fix \(q\geq2\), and put
\[
 G_q=\frac1q\Z/\Z,\qquad U(q)=(\Z/q\Z)^\times.
\]
Let \(\mathcal A\) be a commutative unital \(\C\)-algebra and choose
arbitrary \(t_p\in\mathcal A\), one for each prime \(p\mid q\).
The choices may vanish or be nilpotent. Let
\[
 \mathcal V_q=\bigoplus_{x\in G_q}\mathcal A[x].
\]
For \(p\mid q\) and \(x\in pG_q\), define the prime distribution row
\begin{equation}
 D_p(x)=\sum_{\substack{y\in G_q\\py=x}}[y]-t_p[x].
 \label{eq:dist-prime-row}
\end{equation}
Write \(\mathcal R_q\) for the \(\mathcal A\)-span of these rows and
\(\mathcal U_q=\mathcal V_q/\mathcal R_q\).
The unit group acts by \(u[x]=[ux]\), and each relation map is
equivariant.

For \(d\mid q\), set \(t_d=\prod_{p\mid d}t_p^{v_p(d)}\).
Rows for composite divisors introduce no additional relations. Indeed,
if \(ab\mid q\) and \(x\in abG_q\), then
\begin{equation}
 D_{ab}(x)=
 \sum_{\substack{z\in G_q\\az=x}}D_b(z)+t_bD_a(x).
 \label{eq:dist-composite-row}
\end{equation}
Every \(z\) in this sum belongs to \(bG_q\), so the rows on the right
are legitimate grid rows. Expanding the two sums proves the identity.
Induction on the number of prime factors of \(d\), counted with
multiplicity, proves the assertion for every divisor.

Let \(\chi\) run through all characters of \(U(q)\), and let \(f=f_\chi\)
be the conductor of its primitive inducing character \(\chi^*\).
For \(f\mid d\mid q\), let \(\chi_d\) denote the induced character
modulo \(d\), and define
\begin{equation}
 Y_{d,\chi}=\sum_{a\in U(d)}\chi_d(a)[a/d].
 \label{eq:dist-character-vector}
\end{equation}
For the principal character at \(d=1\), this means
\(Y_{1,\mathbf1}=[0]\). These are formal vectors; no numerical
\(L\)-value is being declared independent.

\begin{lemma}[All character components of the rows]
\label{lem:dist-all-components}
The \(\chi^{-1}\)-eigenspace of \(\mathcal V_q\) is free over
\(\mathcal A\), with basis
\[
 \{Y_{d,\chi}: f_\chi\mid d\mid q\}.
\]
Its relation submodule is generated exactly by the edges
\begin{equation}
 \begin{aligned}
 Y_{pd,\chi}-t_pY_{d,\chi},
     &\qquad p\mid d,\\
 Y_{pd,\chi}-(t_p-\chi^*(p))Y_{d,\chi},
     &\qquad p\nmid d,
 \end{aligned}
 \qquad f_\chi\mid d,\quad pd\mid q.
 \label{eq:dist-character-edges}
\end{equation}
In particular, there are no additional character relations arising
from denominators smaller than, or not divisible by, \(f_\chi\).
\end{lemma}

\begin{proof}
The elements of \(G_q\) of exact denominator \(d\) form a single
\(U(q)\)-orbit, naturally identified with \(U(d)\). Reduction
\(U(q)\to U(d)\) is surjective. The character \(\chi\) factors through
this map exactly when \(f_\chi\mid d\); otherwise its projection on
that orbit is zero. More explicitly, the idempotent
\[
 E_\chi=\frac1{\varphi(q)}\sum_{u\in U(q)}\chi(u)u
\]
satisfies, for \(a\in U(d)\),
\begin{equation}
 E_\chi[a/d]=
 \begin{cases}
 \displaystyle\frac{\chi_d(a)^{-1}}{\varphi(d)}Y_{d,\chi},
       &f_\chi\mid d,\\[6pt]
 0,&f_\chi\nmid d.
 \end{cases}
 \label{eq:dist-projector}
\end{equation}
This follows by grouping \(u\) according to its residue modulo \(d\)
and using character orthogonality on the kernel of reduction.
The idempotents are mutually orthogonal and sum to the identity.
All their scalar denominators are units in \(\mathcal A\), so the
argument applies equally to a nonreduced algebra.

For a fixed \(p\mid q\), consider the equivariant map
\[
 D_p:\mathcal A[pG_q]\longrightarrow\mathcal V_q,\qquad
 [x]\longmapsto D_p(x).
\]
The exact-denominator orbits in its domain have \(pd\mid q\).
By \eqref{eq:dist-projector}, their only nonzero \(\chi\)-components
have \(f_\chi\mid d\), and each such component has the single
generator \(Y_{d,\chi}\). Consequently, computing
\(D_p(Y_{d,\chi})\) computes \emph{every} projected relation for
that \(p\); an orbit on which the projector vanishes cannot create
a hidden relation in the target.

If \(p\mid d\), all \(p\) preimages of each \(a/d\), \(a\in U(d)\),
have exact denominator \(pd\). Summing with weights \(\chi_d(a)\)
therefore gives \(Y_{pd,\chi}-t_pY_{d,\chi}\).
If \(p\nmid d\), one preimage has denominator \(d\), while the other
\(p-1\) have denominator \(pd\). In the former preimage,
\(pc\equiv a\pmod d\), so its character coefficient is
\(\chi_d(a)=\chi^*(p)\chi_d(c)\). Thus the sum is
\[
 Y_{pd,\chi}+\chi^*(p)Y_{d,\chi}-t_pY_{d,\chi}.
 \]
Here \(p\nmid f_\chi\), as \(f_\chi\mid d\), so \(\chi^*(p)\)
is a nonzero root of unity. This proves both the edge formulas
and their completeness.
\end{proof}

\begin{theorem}[Polynomial normal form]
\label{thm:dist-normal-form}
For \(f=f_\chi\mid d\mid q\), define
\begin{equation}
 A_{d,\chi}(\mathbf t)=
 \prod_{p\mid f}t_p^{\,v_p(d)-v_p(f)}
 \prod_{\substack{p\mid d\\p\nmid f}}
       t_p^{\,v_p(d)-1}(t_p-\chi^*(p)).
 \label{eq:dist-normal-coefficient}
\end{equation}
Empty products are \(1\), so \(A_{f,\chi}=1\).
Then the full \(\chi\)-relation submodule is
\begin{equation}
 \mathcal R_{q,\chi}
 =\bigoplus_{\substack{f\mid d\mid q\\d\ne f}}
   \mathcal A\bigl(Y_{d,\chi}-A_{d,\chi}Y_{f,\chi}\bigr).
 \label{eq:dist-triangular-relations}
\end{equation}
In particular,
\begin{equation}
 \mathcal U_q\cong
 \bigoplus_{\chi\in\widehat{U(q)}}\mathcal A\,Y_{f_\chi,\chi}
 \cong\mathcal A^{\varphi(q)}.
 \label{eq:dist-free}
\end{equation}
The character normal form of every vector is obtained by replacing
\(Y_{d,\chi}\) with \(A_{d,\chi}Y_{f_\chi,\chi}\).
\end{theorem}

\begin{proof}
The polynomials \(A_{d,\chi}\) satisfy precisely the edge recursions
\eqref{eq:dist-character-edges}: a prime already dividing \(d\)
contributes \(t_p\), and a prime introduced for the first time
outside \(f\) contributes \(t_p-\chi^*(p)\).
The factors for distinct primes commute.

Follow any chain of prime extensions from \(f\) to \(d\).
Successively eliminating the intermediate \(Y\)'s expresses
\(Y_{d,\chi}-A_{d,\chi}Y_{f,\chi}\) as a combination of the
edge relations. Conversely, write an edge as
\(Y_{pd,\chi}-cY_{d,\chi}\), where
\(A_{pd,\chi}=cA_{d,\chi}\). It equals
\[
 \bigl(Y_{pd,\chi}-A_{pd,\chi}Y_{f,\chi}\bigr)
 -c\bigl(Y_{d,\chi}-A_{d,\chi}Y_{f,\chi}\bigr).
 \]
Lemma~\ref{lem:dist-all-components} therefore gives equality of the
two relation spans.

The displayed generators in \eqref{eq:dist-triangular-relations}
are independent: the coefficient of \(Y_{d,\chi}\), for \(d\ne f\),
in a linear combination is exactly the coefficient of its own
generator. Together with \(Y_{f,\chi}\), they form a triangular
basis of the whole character block. Its quotient is consequently
free of rank one. Summing over all \(\varphi(q)\) characters proves
\eqref{eq:dist-free}.
\end{proof}

\begin{remark}
The proof divides by neither \(t_p\) nor \(t_p-\chi^*(p)\).
It proves freeness first over \(\C[t_p:p\mid q]\) and then over
every \(\C\)-algebra obtained by specialization. Exceptional weights
can therefore make a preferred coordinate vanish, but cannot change
the rank of the full quotient. This distinction is essential when
Euler factors are used to solve a relation system.
\end{remark}

\subsection{Anchors, reflection, and the ranks in the drafts}

In this subsection take \(\mathcal A=\C\).
In a Hurwitz-zeta grid, use the representatives \(0<a\leq1\):
the group element \(0\) represents \(a=1\), not the singular
Hurwitz parameter \(a=0\). Treating the endpoint as known means
imposing the homogeneous anchor \([0]=0\). It removes exactly the
principal generator \(Y_{1,\mathbf1}\). Hence
\begin{equation}
 \dim_\C\mathcal V_q/(\mathcal R_q+\C[0])=\varphi(q)-1.
 \label{eq:dist-anchor-dimension}
\end{equation}
There are \(q-1\) remaining grid variables, so their distribution
coefficient matrix has rank \(q-\varphi(q)\). At rational weights
\(t_p\), the matrix has rational entries and this is also its rank
over \(\Q\).

Let \(J[x]=[-x]\). On the \(\chi\)-block, \(J\) acts as
\(\chi(-1)\). Thus the additional homogeneous reflection rows
\begin{equation}
 [x]+(-1)^k[-x]=0,\qquad x\ne0,
 \label{eq:dist-reflection-rows}
\end{equation}
retain exactly the characters with
\(\chi(-1)=(-1)^{k+1}\). After anchoring, the omission of \(x=0\)
in \eqref{eq:dist-reflection-rows} has no effect.

\begin{corollary}[All-denominator formal ranks]
\label{cor:dist-all-ranks}
For \(q\geq3\) and every integer \(k\geq0\), the formal quotient
by distribution, endpoint, and the reflection rows
\eqref{eq:dist-reflection-rows} has dimension
\begin{equation}
 r_{q,k}=
 \begin{cases}
 \varphi(q)/2-1,&k\ \text{odd},\\
 \varphi(q)/2,&k\ \text{even}.
 \end{cases}
 \label{eq:dist-parity-dimension}
\end{equation}
This conclusion holds for arbitrary prime weights.
In particular it proves the all-\(q\), all-\(k\) formal rank law
for the negative Hurwitz-zeta derivative ladder.
The formal distribution-only rank in the Stieltjes parameter
derivative tower is \(q-\varphi(q)\) for every \(q\geq2\),
every parameter derivative order \(k\geq1\), and every Stieltjes
index \(n\geq1\), relative to the endpoint and lower-index data.
\end{corollary}

\begin{proof}
For \(q\geq3\), the element \(-1\) is a nonidentity involution in
\(U(q)\). Character orthogonality gives
\(\sum_\chi\chi(-1)=0\), and each summand is \(1\) or \(-1\).
There are consequently \(\varphi(q)/2\) characters of each parity.
Reflection retains the parity stated above. The principal
character is even, so the endpoint removes one additional
dimension exactly when \(k\) is odd. This proves
\eqref{eq:dist-parity-dimension}.

For the negative ladder, differentiation of the Hurwitz
multiplication formula at \(s=-k\) gives
\begin{equation}
 \sum_{j=0}^{d-1}
 \zeta'\!\left(-k,\frac{a+j}{d}\right)
 =d^{-k}\bigl(\zeta'(-k,a)+\log d\,\zeta(-k,a)\bigr).
 \label{eq:dist-negative-jet}
\end{equation}
The Bernoulli-polynomial values \(\zeta(-k,a)\) are lower data.
Thus the homogeneous distribution weight is \(t_p=p^{-k}\).
The homogeneous reflection is
\eqref{eq:dist-reflection-rows}; an explicit right-hand side is
given below. This applies the first part to the ladder in the
source article.

For the parameter derivative tower, the Laurent expansion is
\begin{equation}
 \zeta(s,a)=\frac1{s-1}
   +\sum_{n\geq0}\frac{(-1)^n}{n!}\gamma_n(a)(s-1)^n.
 \label{eq:dist-stieltjes-laurent}
\end{equation}
For \(k\geq1\), its pole disappears after differentiation in \(a\):
\begin{equation}
 F_k(\epsilon,a):=
 \partial_a^k\zeta(1+\epsilon,a)
 =(-1)^k(1+\epsilon)_k\zeta(k+1+\epsilon,a)
 =\sum_{n\geq0}\frac{(-1)^n}{n!}
          \gamma_n^{(k)}(a)\epsilon^n.
 \label{eq:dist-stieltjes-generating}
\end{equation}
This standard identity is the starting point of the all-order
parameter derivative formula of Coffey \cite{Coffey}.
Differentiating the Hurwitz multiplication formula \(k\) times
in \(a\), including its chain-rule factors, gives
\[
 \sum_{j=0}^{d-1}
 F_k\!\left(\epsilon,\frac{a+j}{d}\right)
 =d^{k+1+\epsilon}F_k(\epsilon,a).
\]
Comparing coefficients yields the exact triangular identity
\begin{equation}
 \sum_{j=0}^{d-1}
 \gamma_n^{(k)}\!\left(\frac{a+j}{d}\right)
 =d^{k+1}\sum_{r=0}^{n}
   \binom nr(-\log d)^{n-r}\gamma_r^{(k)}(a).
 \label{eq:dist-stieltjes-multiplication}
\end{equation}
With lower indices \(r<n\) treated as specified data, the
homogeneous weight is \(t_p=p^{k+1}\), independently of \(n\).
Equation~\eqref{eq:dist-anchor-dimension} now gives rank
\(q-\varphi(q)\).
\end{proof}

For completeness, the right-hand side of the negative reflection
can be written without introducing another Hurwitz derivative.
For \(0<a<1\), set
\[
 C_j(a)=\sum_{m\geq1}\frac{\cos(2\pi ma)}{m^j},
 \qquad
 S_j(a)=\sum_{m\geq1}\frac{\sin(2\pi ma)}{m^j}.
\]
For \(k\geq1\), these series at \(j=k+1\) converge absolutely, and
the Fourier form of the Hurwitz functional equation \cite{DLMF}
gives
\begin{equation}
 \zeta'(-k,a)+(-1)^k\zeta'(-k,1-a)
 =\frac{(-1)^{\lfloor k/2\rfloor}k!}{(2\pi)^k}
 \begin{cases}
 C_{k+1}(a),&k\ \text{even},\\
 S_{k+1}(a),&k\ \text{odd}.
 \end{cases}
 \label{eq:dist-negative-reflection}
\end{equation}
Indeed, the Fourier expansion of \(\zeta(s,a)\) has the factor
\[
 \frac{2\Gamma(1-s)}{(2\pi)^{1-s}}
 \left(\sin\frac{\pi s}{2}\,C_{1-s}(a)
       +\cos\frac{\pi s}{2}\,S_{1-s}(a)\right).
\]
The indicated reflection selects a sine or cosine prefactor
that vanishes at \(s=-k\). Only the derivative of that prefactor
survives, giving \eqref{eq:dist-negative-reflection}.
This also explains why reflection is an inhomogeneous relation
with depth-one data in the inventory.

\begin{corollary}[The formal gamma grid]
\label{cor:dist-gamma-rank}
For \(q\geq3\), the \(q-1\) formal coordinates
\(\log\Gamma(a/q)\), \(1\leq a<q\), modulo multiplication and
reflection and relative to their stated elementary right-hand
sides, have quotient dimension \(\varphi(q)/2\).
\end{corollary}

\begin{proof}
Lerch's identity
\(\zeta'(0,a)=\log\Gamma(a)-\tfrac12\log(2\pi)\)
translates gamma multiplication into
\eqref{eq:dist-negative-jet} with \(k=0\), hence \(t_p=1\).
The reflection formula
\[
 \zeta'(0,a)+\zeta'(0,1-a)=-\log(2\sin\pi a)
\]
has homogeneous part \([x]+[-x]=0\). Apply
\eqref{eq:dist-parity-dimension} with \(k=0\).
\end{proof}

\begin{remark}[What these rank theorems settle]
\label{rem:dist-formal-only}
The first Stieltjes source states a finite verification followed
by an all-denominator conjecture; the parameter-derivative source
states a general rank theorem while reporting only a finite
calculation. Corollary~\ref{cor:dist-all-ranks} supplies a proof
for the specified relation modules in both cases. The README
had already identified the second proof gap.
Corollary~\ref{cor:dist-gamma-rank} likewise determines the rank
of the stated gamma relation matrix. None of these results says
that the evaluated numbers admit no further relations.
In particular, the claim that multiplication and reflection
generate every algebraic multiplicative relation among rational
gamma values remains part of the conjectural completeness
problem, not a consequence of the matrix rank.
\end{remark}

\subsection{Analytic weights, exceptional coordinates, and finite jets}

Specialize \(t_p=p^s=\exp(s\log p)\), using the real logarithm
of the prime. Formula~\eqref{eq:dist-normal-coefficient} becomes
\begin{equation}
 A_{d,\chi}(s)=
 \left(\frac d{f_\chi}\right)^s
 \prod_{\substack{p\mid d\\p\nmid f_\chi}}
       \left(1-\chi^*(p)p^{-s}\right).
 \label{eq:dist-euler-coefficient}
\end{equation}
This is exactly the familiar imprimitive Euler-factor pattern
for Dirichlet \(L\)-functions \cite{DLMF}. The normal-form proof
shows why it appears in the formal relations before any
evaluation is made.

The symbols of exact denominator \(q\) have character coordinates
\(Y_{q,\chi}\). In the quotient, their \(\chi\)-coordinate is
\begin{equation}
 Y_{q,\chi}=A_{q,\chi}(s)Y_{f_\chi,\chi}.
 \label{eq:dist-top-coordinate}
\end{equation}
They form a basis of the full quotient precisely when every
\(A_{q,\chi}(s)\) is nonzero. In particular this holds when
\(\Repart s\ne0\), since \(\chi^*(p)\) has modulus one.
All nonzero integer weights in the two Stieltjes articles are
therefore free of this coordinate obstruction. On the imaginary
axis the top-denominator coordinates can fail, while the
conductor coordinates of Theorem~\ref{thm:dist-normal-form}
remain a basis.

\begin{theorem}[Exact defect of top-denominator jets]
\label{thm:dist-jet-defect}
Fix \(s_0\in\C\) and \(N\geq0\), and work over
\(\mathcal A_N=\C[\epsilon]/(\epsilon^{N+1})\) with
\(t_p=p^{s_0+\epsilon}\), truncated in this ring.
The full distribution module is free of rank \(\varphi(q)\)
over \(\mathcal A_N\), hence has complex dimension
\((N+1)\varphi(q)\).
Set
\begin{equation}
 r_\chi(s_0)=
 \#\{p\mid q:p\nmid f_\chi,\ p^{s_0}=\chi^*(p)\}.
 \label{eq:dist-resonance-count}
\end{equation}
The image of the top-denominator submodule in its \(\chi\)-block
has complex dimension
\begin{equation}
 \max\{N+1-r_\chi(s_0),0\},
 \label{eq:dist-jet-image}
\end{equation}
and its cokernel has dimension
\(\min\{r_\chi(s_0),N+1\}\).
An endpoint anchor or a fixed-parity reflection restricts these
formulas to the retained characters. The resulting full quotient
remains free over \(\mathcal A_N\) with the ranks already computed.
\end{theorem}

\begin{proof}
Freeness is Theorem~\ref{thm:dist-normal-form}, applied to
\(\mathcal A_N\). A resonant factor in
\eqref{eq:dist-euler-coefficient} has a simple zero at \(s_0\):
\[
 \left.\frac{d}{ds}
 (1-\chi^*(p)p^{-s})\right|_{s=s_0}=\log p\ne0.
\]
All other factors, including \((q/f_\chi)^s\), are nonzero
there. In the ring of analytic germs,
\[
 A_{q,\chi}(s_0+\epsilon)
   =\epsilon^{r_\chi(s_0)}u_\chi(\epsilon),
 \qquad u_\chi(0)\ne0.
\]
After truncation, \(u_\chi\) is a unit. The
top-denominator image in the free rank-one block is therefore
the ideal \(\epsilon^{r_\chi(s_0)}\mathcal A_N\).
Its basis consists of the powers
\(\epsilon^{r_\chi},\ldots,\epsilon^N\) when \(r_\chi\leq N\),
and it is zero otherwise. This proves both dimensions.
Anchoring and reflection remove entire character blocks,
so the same calculation applies to each retained block.
\end{proof}

Comparing Taylor coefficients of a weighted distribution identity
gives the corresponding stacked finite-jet system. If one identifies
the coefficient at order \(j\) with
\(\epsilon^{N-j}[x]\), its relation space is the
\(\mathcal A_N\)-span above; factorial normalizations merely
rescale coordinates. Thus freeness also gives the rank of the
complete simultaneous jet system, rather than just its successive
highest layers.

\begin{proposition}[Multiplicity of nonzero resonances]
\label{prop:dist-resonance-multiplicity}
Resonances occur only on \(\Repart s_0=0\).
For \(s_0\ne0\), each \(r_\chi(s_0)\) is at most one.
Multiple resonant primes in a single character block can occur
only at \(s_0=0\).
\end{proposition}

\begin{proof}
Taking absolute values in \(p^{s_0}=\chi^*(p)\) gives
\(\Repart s_0=0\). Write \(s_0=\ii t\), with \(t\ne0\).
Because every character value is a root of unity, resonance
at \(p\) implies
\(t\log p/(2\pi)\in\Q\). This rational number is nonzero.
If distinct primes \(p\) and \(\ell\) both resonate, then
\(\log p/\log\ell\in\Q\). Writing this positive rational as
\(a/b\) would give \(p^b=\ell^a\) for positive integers \(a,b\),
contradicting unique factorization.
\end{proof}

For a concrete origin resonance, take \(q=14\) and let \(\chi^*\)
be the quadratic character modulo \(7\). Its conductor is \(7\),
its parity is odd, and \(\chi^*(2)=1\). Then
\begin{equation}
 Y_{14,\chi}(s)=(2^s-1)Y_{7,\chi}(s).
 \label{eq:dist-resonance-example}
\end{equation}
At \(s=0\), the top-denominator coordinate vanishes although the
conductor direction persists. For any analytic realization of
this relation,
\begin{align}
 Y_{14,\chi}'(0)&=\log2\,Y_{7,\chi}(0),\notag\\
 Y_{14,\chi}''(0)&=
 2\log2\,Y_{7,\chi}'(0)+(\log2)^2Y_{7,\chi}(0).
 \label{eq:dist-resonance-example-jets}
\end{align}
Thus recovering the first conductor derivative from these
top-denominator data requires the second top-denominator jet.
The actual realization
\[
 Y_{d,\chi}(s)
 =\sum_{a\in U(d)}\chi_d(a)\zeta(s,a/d)
 =d^sL(s,\chi_d)
\]
obeys the same identity. This example concerns an imprimitive
Euler factor; it is distinct from the trivial-zero regularization
of the primitive \(L\)-function in the next section.

\section{One normalized functional equation for all derivative bridges}
\label{sec:bridges}

The Stieltjes drafts connect derivatives of \(L(s,\chi)\) at a positive
integer with derivatives at its reflected negative integer. A single
identity of analytic germs organizes all these formulas and treats
trivial zeros without differentiating singular logarithms. It also
makes the conjugation required for complex characters unavoidable.
The input is the classical primitive Dirichlet functional equation
and the reflection and duplication formulas for \(\Gamma\)
\cite{DLMF}; the all-order normalized form below is a systematic
continuation of the bridges used in the corpus.

\subsection{Regularization and the normalized germ}

Let \(\chi\) be primitive of conductor \(q\), allowing the trivial
character of conductor \(1\), for which \(L(s,\chi)=\zeta(s)\).
Fix \(k\geq1\), and write
\[
 \delta=\frac{1-\chi(-1)}2\in\{0,1\},\qquad
 \nu=
 \begin{cases}
 1,&\chi(-1)=(-1)^k,\\
 0,&\chi(-1)\ne(-1)^k.
 \end{cases}
\]
Define
\begin{equation}
 B(t)=\frac{L(-k+t,\chi)}{t^\nu},
 \qquad
 B^\#(z)=\overline{B(\overline z)},
 \label{eq:bridge-regularization}
\end{equation}
with the removable value filled in when \(\nu=1\).
The function \(B^\#\) is holomorphic near zero; it is coefficientwise
conjugation of the Taylor series of \(B\).

\begin{theorem}[Normalized positive--negative germ]
\label{thm:bridge-normalized}
The quantities \(L(k+1,\chi)\) and \(B(0)\) are nonzero. On a
sufficiently small disk about \(z=0\),
\begin{equation}
 \frac{L(k+1+z,\chi)}{L(k+1,\chi)}
 =
 K_\nu(z)\frac{B^\#(-z)}{B^\#(0)},
 \label{eq:bridge-normalized-germ}
\end{equation}
where
\begin{equation}
 \begin{aligned}
 K_\nu(z)&=
 \left(\frac{2\pi}{q}\right)^z
 \frac{\Gamma(k+1)}{\Gamma(k+1+z)}\,T_\nu(z),\\
 T_0(z)&=\sec\frac{\pi z}{2},\qquad
 T_1(z)=\frac{\pi z/2}{\sin(\pi z/2)},\quad T_1(0)=1.
 \end{aligned}
 \label{eq:bridge-elementary-kernel}
\end{equation}
All three normalized factors in
\eqref{eq:bridge-normalized-germ} equal \(1\) at the origin.
\end{theorem}

\begin{proof}
The Euler product gives \(L(k+1,\chi)\ne0\), since \(k+1>1\).
Use the completed function
\[
 \Lambda(s,\chi)=
 \left(\frac q\pi\right)^{(s+\delta)/2}
 \Gamma\!\left(\frac{s+\delta}{2}\right)L(s,\chi)
\]
and its functional equation
\begin{equation}
 \Lambda(s,\chi)=\varepsilon_\chi\Lambda(1-s,\overline\chi),
 \qquad |\varepsilon_\chi|=1.
 \label{eq:bridge-completed-fe}
\end{equation}
At \(s=-k\), the gamma factor is finite and nonzero if
\(\nu=0\), and has a simple pole if \(\nu=1\).
The reflected completed value at \(k+1\) is finite and nonzero.
Consequently \(L(-k,\chi)\ne0\) in the first case, and
\(L(s,\chi)\) has a simple zero at \(-k\) in the second.
In either case \(B\) is holomorphic and \(B(0)\ne0\).

Put
\[
 a=\frac{\delta-k}{2},\qquad
 b=\frac{k+1+\delta}{2}.
\]
Solving \eqref{eq:bridge-completed-fe} at \(s=k+1+z\) gives
\begin{equation}
 L(k+1+z,\chi)
 =\varepsilon_\chi
 \left(\frac q\pi\right)^{-k-\frac12-z}
 \frac{\Gamma(a-z/2)}{\Gamma(b+z/2)}
 L(-k-z,\overline\chi).
 \label{eq:bridge-fe-before-normalization}
\end{equation}
Since
\(L(-k-z,\overline\chi)=(-z)^\nu B^\#(-z)\),
the numerator
\[
 F_\nu(z)=(-z)^\nu\Gamma(a-z/2)
\]
has a nonzero removable value at the origin.
Dividing \eqref{eq:bridge-fe-before-normalization} by its value
at \(z=0\) cancels the root number and all constant factors:
\begin{equation}
 \frac{L(k+1+z,\chi)}{L(k+1,\chi)}
 =
 \left(\frac q\pi\right)^{-z}
 \frac{F_\nu(z)}{F_\nu(0)}
 \frac{\Gamma(b)}{\Gamma(b+z/2)}
 \frac{B^\#(-z)}{B^\#(0)}.
 \label{eq:bridge-normalized-intermediate}
\end{equation}

If \(\nu=0\), then \(a\) is a half-integer.
The gamma reflection formula and
\(\sin(\pi(a-z/2))=\sin(\pi a)\cos(\pi z/2)\)
give
\begin{equation}
 \frac{F_0(z)}{F_0(0)}
 =\frac{\Gamma(1-a)}{\Gamma(1-a+z/2)}
       \sec\frac{\pi z}{2}.
 \label{eq:bridge-gamma-half-integer}
\end{equation}
If \(\nu=1\), write \(a=-m\), \(m\geq0\).
Then \(F_1(0)=2(-1)^m/m!\), and the same reflection formula
gives
\begin{equation}
 \frac{F_1(z)}{F_1(0)}
 =\frac{\Gamma(1-a)}{\Gamma(1-a+z/2)}
       \frac{\pi z/2}{\sin(\pi z/2)}.
 \label{eq:bridge-gamma-integer}
\end{equation}
For example, the pole in \(\Gamma(-m-z/2)\) is canceled by
the explicit factor \(-z\), so no logarithm of a function
vanishing at zero has been taken.

The unordered pair \(\{b,1-a\}\) is
\(\{(k+1)/2,(k+2)/2\}\), for either value of \(\delta\).
Duplication therefore yields
\begin{equation}
 \frac{\Gamma(b)\Gamma(1-a)}
 {\Gamma(b+z/2)\Gamma(1-a+z/2)}
 =2^z\frac{\Gamma(k+1)}{\Gamma(k+1+z)}.
 \label{eq:bridge-gamma-duplication}
\end{equation}
Substituting \eqref{eq:bridge-gamma-half-integer} or
\eqref{eq:bridge-gamma-integer}, followed by
\eqref{eq:bridge-gamma-duplication}, into
\eqref{eq:bridge-normalized-intermediate} proves the theorem.
\end{proof}

\begin{remark}
The statement concerns analytic germs and local logarithms.
It requires no claim about a global zero-free region.
For a complex character the functional equation involves
\(\overline\chi\), and therefore the reflected Taylor coefficients
are conjugated. Replacing \(B^\#\) by \(B\) is justified for a real
character but generally false for a complex one.
\end{remark}

\subsection{Closed logarithmic coefficients at every order}

Choose the normalized local logarithms, each zero at the origin,
of \(L(k+1+z,\chi)/L(k+1,\chi)\) and \(B(t)/B(0)\).
For \(r\geq1\), set
\[
 \lambda_r=
 \left.\frac{d^r}{ds^r}\log L(s,\chi)\right|_{s=k+1},
 \qquad
 b_r=
 \left.\frac{d^r}{dt^r}\log B(t)\right|_{t=0}.
\]
These derivatives do not depend on a choice of logarithm branch.
Write
\[
 H_k^{(r)}=\sum_{j=1}^{k}j^{-r},\qquad H_k=H_k^{(1)},
 \qquad
 \eta(r)=(1-2^{1-r})\zeta(r)\quad(r\geq2).
\]

\begin{theorem}[All-order logarithmic bridges]
\label{thm:bridge-log-jets}
For every \(r\geq1\),
\begin{equation}
 \lambda_r-(-1)^r\overline{b_r}=\kappa_r,
 \label{eq:bridge-log-jets}
\end{equation}
where
\begin{equation}
 \kappa_r=
 \begin{cases}
 \displaystyle\gamma+\log\frac{2\pi}{q}-H_k,
       &r=1,\\[4pt]
 \displaystyle(r-1)!\bigl(\zeta(r)-H_k^{(r)}\bigr),
       &r\geq3\ \text{odd},\\[4pt]
 \displaystyle(r-1)!\bigl(H_k^{(r)}+(-1)^\nu\eta(r)\bigr),
       &r\geq2\ \text{even}.
 \end{cases}
 \label{eq:bridge-kappa}
\end{equation}
\end{theorem}

\begin{proof}
Taking normalized logarithms of
\eqref{eq:bridge-normalized-germ} gives
\[
 \log\frac{L(k+1+z,\chi)}{L(k+1,\chi)}
 =\log K_\nu(z)
  +\overline{\log\frac{B(-\overline z)}{B(0)}}.
\]
The \(r\)-th derivative of the last term at zero is
\((-1)^r\overline{b_r}\). It remains to compute the
derivatives of \(\log K_\nu\).

The convergent gamma expansion, on \(|z|<1\), is
\begin{equation}
 \log\frac{\Gamma(k+1+z)}{\Gamma(k+1)}
 =(H_k-\gamma)z+
 \sum_{r\geq2}\frac{(-1)^r}{r}
       \bigl(\zeta(r)-H_k^{(r)}\bigr)z^r.
 \label{eq:bridge-log-gamma-series}
\end{equation}
The sine and cosine products give, in the same disk,
\begin{align}
 \log\sec\frac{\pi z}{2}
 &=\sum_{m\geq1}
   \frac{(1-2^{-2m})\zeta(2m)}{m}z^{2m},
 \label{eq:bridge-log-sec-series}\\
 \log\frac{\pi z/2}{\sin(\pi z/2)}
 &=\sum_{m\geq1}
   \frac{2^{-2m}\zeta(2m)}{m}z^{2m}.
 \label{eq:bridge-log-sinc-series}
\end{align}
These identities also follow directly by taking logarithms
of the absolutely locally convergent products
\(\cos(\pi z/2)=\prod_{j\geq0}(1-z^2/(2j+1)^2)\) and
\(\sin(\pi z/2)/(\pi z/2)
=\prod_{j\geq1}(1-z^2/(2j)^2)\).

The linear term of \(\log K_\nu\) is the stated
\(\kappa_1\). Odd coefficients of order at least three come
only from the negative of
\eqref{eq:bridge-log-gamma-series}, giving the second case.
For \(r=2m\), the secant coefficient combines with the gamma
coefficient to give
\[
 (r-1)!\left[
  H_k^{(r)}-\zeta(r)+2(1-2^{-r})\zeta(r)\right]
 =(r-1)!\bigl(H_k^{(r)}+\eta(r)\bigr).
\]
Replacing the secant factor by the sine quotient instead gives
\[
 (r-1)!\left[
  H_k^{(r)}-\zeta(r)+2^{1-r}\zeta(r)\right]
 =(r-1)!\bigl(H_k^{(r)}-\eta(r)\bigr).
\]
These are exactly the two even cases in
\eqref{eq:bridge-kappa}.
\end{proof}

At a nonzero reflected value, \(\nu=0\), the first two identities
read
\begin{align}
 \frac{L'(k+1,\chi)}{L(k+1,\chi)}
 +\overline{\frac{L'(-k,\chi)}{L(-k,\chi)}}
 &=\gamma+\log\frac{2\pi}{q}-H_k,
 \label{eq:bridge-first-nonzero}\\
 (\log L)''(k+1,\chi)
 -\overline{(\log L)''(-k,\chi)}
 &=H_k^{(2)}+\frac{\pi^2}{12}.
 \label{eq:bridge-second-nonzero}
\end{align}
At a trivial zero, \(\nu=1\), the Taylor expansion
\[
 B(t)=L'(-k,\chi)
   +\frac{L''(-k,\chi)}2t
   +\frac{L'''(-k,\chi)}6t^2+\cdots
\]
instead gives
\begin{align}
 \frac{L'(k+1,\chi)}{L(k+1,\chi)}
 +\frac12\overline{\frac{L''(-k,\chi)}{L'(-k,\chi)}}
 &=\gamma+\log\frac{2\pi}{q}-H_k,
 \label{eq:bridge-first-zero}\\
 (\log L)''(k+1,\chi)
 -\overline{\left(
    \frac{L'''(-k,\chi)}{3L'(-k,\chi)}
   -\frac{L''(-k,\chi)^2}{4L'(-k,\chi)^2}
   \right)}
 &=H_k^{(2)}-\frac{\pi^2}{12}.
 \label{eq:bridge-second-zero}
\end{align}
These formulas preserve the familiar first-order harmonic
bridge and supply a uniform extension to every derivative
order and both character parities.

\subsection{Raw jets and an exact triangular implementation}

The normalized germ also gives direct formulas for ordinary
derivatives, without first converting each order to logarithmic
derivatives. Define \(K_m\) by
\[
 K_\nu(z)=\sum_{m\geq0}K_m\frac{z^m}{m!},\qquad K_0=1.
\]
Since \(\log K_\nu(z)=\sum_{r\geq1}\kappa_rz^r/r!\),
these coefficients are the complete exponential Bell polynomials:
\begin{equation}
 K_m=
 m!\!\!\sum_{\substack{m_1,\ldots,m_m\geq0\\
                     m_1+2m_2+\cdots+mm_m=m}}
 \prod_{r=1}^{m}\frac1{m_r!}
       \left(\frac{\kappa_r}{r!}\right)^{m_r}
 \qquad(m\geq1).
 \label{eq:bridge-bell-coefficients}
\end{equation}
Leibniz's rule in \eqref{eq:bridge-normalized-germ} proves
\begin{equation}
 \frac{L^{(m)}(k+1,\chi)}{L(k+1,\chi)}
 =
 \sum_{j=0}^{m}\binom mj K_{m-j}(-1)^j
       \overline{\frac{B^{(j)}(0)}{B(0)}}.
 \label{eq:bridge-raw-jets}
\end{equation}
The required regularized derivatives are explicitly
\begin{equation}
 \frac{B^{(j)}(0)}{B(0)}
 =\frac{j!\,\nu!}{(j+\nu)!}
    \frac{L^{(j+\nu)}(-k,\chi)}{L^{(\nu)}(-k,\chi)}.
 \label{eq:bridge-regularized-raw-jets}
\end{equation}
The diagonal coefficient in
\eqref{eq:bridge-raw-jets} is \((-1)^m\), so the transformation
is triangular and invertible. Its coefficients belong to the
ring generated by the explicit \(\kappa_r\)'s. This supplies
exact relation certificates at arbitrary order and avoids
repeated symbolic differentiation across gamma poles.

\begin{remark}[Correction of the complex-character formulas]
The older derivative-relations report omits the conjugation
in versions of the trivial-zero first bridge and the nonzero
second bridge. Equations~\eqref{eq:bridge-first-zero} and
\eqref{eq:bridge-second-nonzero} give the required replacements.
The two main Stieltjes articles already include conjugation
in their displayed first bridges; those correct formulas
should be retained. Exact source locators and numerical
counterexamples to the omitted-conjugation versions appear
in the correction register.
\end{remark}

\section{High parameter derivatives of Stieltjes constants}
\label{sec:asymptotics}

The parameter-derivative tower admits a useful separation into a
finite polynomial coefficient and an exponentially small spectral
tail. This gives a different asymptotic regime from the frequently
studied limit of a large Stieltjes index with derivative order fixed.
Here the derivative order \(k\) tends to infinity, while the index
\(n\) may grow proportionally to \(\log k\). The resulting profile is
the reciprocal gamma function.

\subsection{An exact spectral expansion}

Fix the Laurent convention
\begin{equation}\label{eq:stieltjes-convention}
\zeta(s,a)=\frac{1}{s-1}
 +\sum_{n=0}^{\infty}\frac{(-1)^n}{n!}\gamma_n(a)(s-1)^n,
\qquad a>0.
\end{equation}
The superscript \(\gamma_n^{(k)}(a)\) always denotes differentiation
with respect to \(a\), not to the index or the zeta variable. Define
\begin{equation}\label{eq:coefficient-polynomial}
P_{n,k}(x)=n![z^n]\left(e^{xz}\prod_{j=1}^{k}
                  \left(1+\frac zj\right)\right).
\end{equation}
This is a degree-\(n\) polynomial in \(x\); its coefficients are
rational. In particular \(P_{0,k}=1\) and
\(P_{1,k}(x)=x+H_k\).

\begin{theorem}[Spectral expansion and a finite-tail certificate]
\label{thm:spectral}
For \(a>0\), integers \(k\ge1\), and \(n\ge0\),
\begin{equation}\label{eq:spectral}
\frac{(-1)^{n+k}}{k!}\gamma_n^{(k)}(a)
=\sum_{m=0}^{\infty}
 \frac{P_{n,k}(-\log(a+m))}{(a+m)^{k+1}}.
\end{equation}
The series converges absolutely, locally uniformly for \(a>0\).
If the sum is truncated before \(m=M\), where \(b=a+M\ge1\),
then for every real \(0<R<k\) the magnitude of the omitted part is at
most
\begin{equation}\label{eq:spectral-tail}
\frac{n!}{R^n}\prod_{j=1}^{k}\left(1+\frac Rj\right)
\left(b^{-k-1+R}+\frac{b^{-k+R}}{k-R}\right).
\end{equation}
\end{theorem}

\begin{proof}
The Hurwitz-zeta parameter identity
\[
\partial_a^k\zeta(s,a)=(-1)^k(s)_k\zeta(s+k,a)
\]
follows initially by differentiating its absolutely convergent
Dirichlet series and then by analytic continuation; here
\((s)_k=s(s+1)\cdots(s+k-1)\).
The pole in \eqref{eq:stieltjes-convention} is independent of \(a\).
Put \(s=1+z\), use
\((1+z)_k=k!\prod_{j=1}^k(1+z/j)\), and compare coefficients.
For \(|z|<k\), the series for \(\zeta(k+1+z,a)\) converges normally.
This gives \eqref{eq:spectral} and justifies the coefficient
interchange. Alternatively, absolute convergence follows directly
because each summand is a linear combination of
\((a+m)^{-k-1}\log^\ell(a+m)\), \(0\le\ell\le n\).
The parameter identity is also the starting point of
Coffey's all-order formulas~\cite{Coffey}.

On the circle \(|z|=R\), and for \(x\ge b\ge1\),
\[
\left|e^{-z\log x}\prod_{j=1}^k(1+z/j)\right|
\le x^R\prod_{j=1}^k(1+R/j).
\]
Cauchy's coefficient estimate bounds the \(m\)-th omitted summand
by \(n!R^{-n}\prod_j(1+R/j)(a+m)^{-k-1+R}\).
For the positive decreasing power, comparison with its integral gives
\[
\sum_{m=M}^{\infty}(a+m)^{-k-1+R}
\le b^{-k-1+R}+\int_b^\infty x^{-k-1+R}\,dx.
\]
This is \eqref{eq:spectral-tail}.
\end{proof}

The condition \(b\ge1\) belongs to the stated majorant, rather than to
the exact expansion. It always holds for the tail after the first
term, because then \(b=a+1>1\). The choice of \(R\) can be optimized
for a given \(n,k,M\); there is no need to differentiate a
highly oscillatory numerical representation of \(\gamma_n\).

\subsection{Exact rational enclosures at parameter one}

Write \(\stirling{v}{u}\) for the unsigned Stirling number of the
first kind. The elementary identity
\begin{equation}\label{eq:stirling-product}
\prod_{j=1}^k(1+z/j)
 =\frac1{k!}\sum_{n=0}^k\stirling{k+1}{n+1}z^n
\end{equation}
turns the first spectral term at \(a=1\) into an exact rational
number.

\begin{corollary}[Rational certificates]\label{cor:rational}
Let \(k\ge2\), \(n\ge0\), and
\[
e_n(k)=
\begin{cases}
\stirling{k+1}{n+1}/k!,&0\le n\le k,\\
0,&n>k.
\end{cases}
\]
For any integer \(1\le h<k\), set
\begin{equation}\label{eq:rational-radius}
B_{k,n,h}
=h^{-n}\binom{k+h}{h}
 \left(2^{-k-1+h}+\frac{2^{-k+h}}{k-h}\right).
\end{equation}
Then the following is an exact inclusion between real numbers:
\begin{equation}\label{eq:rational-inclusion}
\frac{(-1)^{n+k}\gamma_n^{(k)}(1)}{k!\,n!}
\in[e_n(k)-B_{k,n,h},\,e_n(k)+B_{k,n,h}].
\end{equation}
One may take the minimum radius over all \(h=1,\ldots,k-1\).
In particular, the simple choice \(h=1\) gives
\begin{equation}\label{eq:simple-rational-radius}
B_k=(k+1)\left(2^{-k}+\frac{2^{-k+1}}{k-1}\right),
\end{equation}
independently of \(n\).
\end{corollary}

\begin{proof}
Apply Theorem~\ref{thm:spectral} with \(a=M=1\), \(b=2\), and
\(R=h\), and divide its bound by \(n!\). The product in that bound is
\(\prod_{j=1}^k(1+h/j)=\binom{k+h}{h}\).
Equation~\eqref{eq:stirling-product} supplies the center.
All quantities in the resulting endpoints are rational.
\end{proof}

Thus the term \emph{certificate} has its literal meaning here:
integer arithmetic constructs the endpoints, and the proved tail
bound supplies the inclusion. Floating-point agreement is not being
renamed a proof. For example, positivity of an endpoint
\(e_n(k)-B_{k,n,h}\) rigorously proves
\((-1)^{n+k}\gamma_n^{(k)}(1)>0\).
The accompanying script emits centers, radii, and both endpoints
as numerator--denominator pairs, and checks their arithmetic exactly.

\subsection{A uniform coefficient lemma}

The main asymptotic step is an elementary coefficient estimate.
It is convenient to set
\(\fall{n}{j}=n(n-1)\cdots(n-j+1)\) for \(j\le n\), and
\(\fall{n}{j}=0\) for \(j>n\).

\begin{lemma}[Two uniform coefficient terms]\label{lem:coefficient}
Suppose \(f\) is analytic on a neighborhood of the closed disk
\(|z|\le R\), with \(|f(z)|\le M\) there. Let \(L>0\),
\(n\ge1\), and \(r=n/L<R\). Put
\[
Q_n(L;f)=\frac{n!}{L^n}[z^n]e^{Lz}f(z).
\]
Then
\begin{align}
|Q_n(L;f)-f(r)|
&\le\frac{Mn}{L^2R^2(1-r/R)^3},\label{eq:coefficient-first}\\
\left|Q_n(L;f)-f(r)+\frac{n}{2L^2}f''(r)\right|
&\le\frac{3Mn}{L^3R^3(1-r/R)^5}.\label{eq:coefficient-second}
\end{align}
For \(n=0\), \(Q_0(L;f)=f(0)\) exactly.
\end{lemma}

\begin{proof}
Write \(f(z)=\sum_{j\ge0}c_jz^j\). Since \(|c_j|\le MR^{-j}\),
\[
Q_n(L;f)=\sum_{j\ge0}c_j\frac{\fall{n}{j}}{L^j}
\]
is a finite sum, which we compare with the absolutely convergent
Taylor series at \(r\).

For completeness, the needed inequalities remain valid when \(j>n\).
Choose \(j\) independent uniform elements of a set of \(n\) elements.
The probability of no repeated choice is
\(\fall{n}{j}/n^j\), including its value zero when \(j>n\).
For the \(K=\binom j2\) pair-collision events, every event has
probability \(1/n\), and every intersection of two distinct events
has probability \(1/n^2\), even if the pairs share a position.
The union bound and the first two Bonferroni inequalities give
\begin{align*}
0&\le n^j-\fall{n}{j}\le\binom j2 n^{j-1},\\
0&\le\fall{n}{j}-n^j+\binom j2 n^{j-1}
 \le\binom{\binom j2}{2}n^{j-2}.
\end{align*}
At \(j=0,1,2\), the second remainder is zero, so no negative-power
exception is needed. Now use
\[
\sum_{j\ge2}\binom j2 x^{j-2}=(1-x)^{-3},
\qquad
\sum_{j\ge3}\binom{\binom j2}{2}x^{j-3}=3(1-x)^{-5}.
\]
The first inequality gives \eqref{eq:coefficient-first}.
In the second comparison,
\[
\sum_{j\ge2}c_j\binom j2\frac{n^{j-1}}{L^j}
=\frac{n}{2L^2}f''(r).
\]
This proves \eqref{eq:coefficient-second}.
\end{proof}

\subsection{The reciprocal-gamma profile}

\begin{theorem}[A joint logarithmic-range asymptotic]\label{thm:profile}
Fix \(0<a_0\le a_1<\infty\) and \(C>0\). For
\(a\in[a_0,a_1]\), put
\[
L=\log(k/a),\qquad r=\frac nL,\qquad
g(r)=\frac1{\Gamma(1+r)},\qquad
\rho=\frac{a_1}{a_1+1}<1.
\]
Uniformly over integers \(0\le n\le CL\), as \(k\to\infty\),
\begin{equation}\label{eq:uniform-profile}
\frac{(-1)^{n+k}a^{k+1}\gamma_n^{(k)}(a)}{k!\,L^n}
=g(r)-\frac{n}{2L^2}g''(r)
 +O_{a_0,a_1,C}\left(\frac n{L^3}+k^{C+1}\rho^k\right).
\end{equation}
Equivalently, the two displayed terms are
\begin{equation}\label{eq:profile-polygamma}
\frac1{\Gamma(1+r)}
\left[1+\frac{n}{2L^2}
 \left\{\psi^{(1)}(1+r)-\psi(1+r)^2\right\}\right].
\end{equation}
The formula includes \(n=0\), for which the first spectral term is
exactly one. In particular, if \(n/\log(k/a)\to c\ge0\) in a bounded
range, then the left side of \eqref{eq:uniform-profile} tends to
\(1/\Gamma(1+c)\).
\end{theorem}

\begin{proof}
Separate the \(m=0\) term of \eqref{eq:spectral}. Its generating
function factors as
\[
e^{-z\log a}\prod_{j=1}^k(1+z/j)=e^{Lz}f_k(z),
\qquad
f_k(z)=k^{-z}\prod_{j=1}^k(1+z/j).
\]
The gamma product gives
\[
f_k(z)=\frac{\Gamma(k+1+z)}
 {\Gamma(k+1)k^z\Gamma(1+z)}
      =g(z)e^{D_k(z)}.
\]
On \(|z|<k+1\), define the analytic function
\begin{equation}\label{eq:Dk}
D_k(z)=z(H_k-\log k-\gamma)
 +\sum_{\ell=2}^{\infty}
   \frac{(-1)^\ell z^\ell}{\ell}
       \bigl(\zeta(\ell)-H_k^{(\ell)}\bigr).
\end{equation}
The factorization follows first on a small disk from the
Weierstrass product and then throughout \(|z|<k+1\) by analytic
continuation. Defining \(D_k\) by its series avoids taking the
logarithm of \(g\) at its zeros.

The elementary bounds
\[
0<H_k-\log k-\gamma<\frac1{2k},\qquad
\sum_{j>k}j^{-2}\le\frac1k
\]
imply, for \(|z|\le R<k+1\),
\begin{equation}\label{eq:Dk-bound}
|D_k(z)|\le\frac{R}{2k}
   +\frac{R^2}{2k(1-R/(k+1))}.
\end{equation}
Indeed, for \(\ell\ge2\), bound the tail of \(\zeta(\ell)\) by
\((k+1)^{2-\ell}/k\), and use \(1/\ell\le1/2\).
Consequently \(f_k\) is uniformly bounded on every fixed disk,
\(f_k-g=O_R(k^{-1})\), and Cauchy's estimate gives
\(f_k''-g''=O_C(k^{-1})\) on \([0,C]\).
The same estimate with \(R=r\), or the real product, also gives the
slightly sharper
\[
f_k(r)-g(r)=O_C(r/k)\qquad(0\le r\le C).
\]
It is this factor \(r\) that preserves the exact \(n=0\) first term.

Choose \(R=C+1\) in Lemma~\ref{lem:coefficient}. For \(n\ge1\)
the normalized first spectral term is
\[
Q_n(L;f_k)
=f_k(r)-\frac{n}{2L^2}f_k''(r)+O_C(n/L^3).
\]
Replacing \(f_k\) by \(g\) costs
\(O_C(n/(kL)+n/(kL^2))\), which is absorbed in \(O(n/L^3)\)
because \(L^2\le k\) for all sufficiently large \(k\), uniformly
on the compact parameter interval.

It remains to control the other spectral terms. Apply
\eqref{eq:spectral-tail} with \(M=1\), and multiply the bound by
\(a^{k+1}/L^n\). The result is
\[
\frac{n!}{(RL)^n}\prod_{j=1}^k(1+R/j)
\left(\frac{a}{a+1}\right)^{k+1}(a+1)^R
\left(1+\frac{a+1}{k-R}\right).
\]
The first factor is at most one: for \(n\ge1\), use
\(n!\le n^n\) and \(n/L\le C<R\); for \(n=0\), it is one by
convention. The product is \(O_R(k^R)\), as follows already from
\eqref{eq:Dk-bound} at \(z=R\).
The remaining factors are uniformly bounded apart from
\((a/(a+1))^{k+1}\le\rho^{k+1}\). Thus the normalized tail is
\(O_{a_0,a_1,C}(k^{C+1}\rho^k)\). The case \(n=0\) follows
directly from \(P_{0,k}=1\).
Finally,
\(g''(r)=g(r)\{\psi(1+r)^2-\psi^{(1)}(1+r)\}\).
\end{proof}

\begin{corollary}[A uniform eventual sign]\label{cor:sign}
For every compact interval \([a_0,a_1]\subset(0,\infty)\) and
every \(C>0\), there is \(k_0\) such that
\[
(-1)^{n+k}\gamma_n^{(k)}(a)>0
\]
for all \(k\ge k_0\), \(a\in[a_0,a_1]\), and integers
\(0\le n\le C\log(k/a)\).
\end{corollary}

\begin{proof}
The continuous function \(g\) is strictly positive on \([0,C]\),
so its minimum there is positive. Both the displayed correction
and the remainder in Theorem~\ref{thm:profile} tend uniformly to zero.
The other factors in the normalization are positive.
\end{proof}

The corollary is restricted to its stated range. It gives no sign
theorem when \(n\) grows arbitrarily quickly compared with \(\log k\).
For a concrete finite pair \(n,k\) at \(a=1\), the rational interval
in Corollary~\ref{cor:rational} is a separate, effective certificate.

\begin{figure}[tbp]
\centering
\includegraphics[width=\textwidth]{figures/derivative-profile.pdf}
\caption{The reciprocal-gamma profile at \(a=1\).
The first spectral term is computed from the elementary-symmetric
coefficient recurrence; its rigorously bounded spectral error is
smaller than the plotted resolution.
The lower panel compares its deviation from the limiting profile
with its deviation after the explicit second term is included.
These plotted arithmetic evaluations are floating point; the separate
rational certificates do not depend on them.}
\label{fig:profile}
\end{figure}

\subsection{Fixed-index expansions and the classical mechanism}

For a fixed \(n\), it is also useful to retain the polynomial in
\(L=\log(k/a)\). Define Appell polynomials
\[
A_n(L)=n![z^n]\frac{e^{Lz}}{\Gamma(1+z)}.
\]
The first few are
\begin{align*}
A_0(L)&=1,\\
A_1(L)&=L+\gamma,\\
A_2(L)&=(L+\gamma)^2-\zeta(2),\\
A_3(L)&=(L+\gamma)^3-3\zeta(2)(L+\gamma)+2\zeta(3).
\end{align*}
The standard gamma-ratio expansion~\cite{DLMF} gives, uniformly on
fixed \(z\)-disks,
\[
e^{D_k(z)}=
1+\frac{z+z^2}{2k}
 +\frac{3z^4+2z^3-3z^2-2z}{24k^2}+O(k^{-3}).
\]
Consequently, for fixed \(n\ge1\) and \(a\) in a compact positive
interval,
\begin{align}
\frac{(-1)^{n+k}a^{k+1}\gamma_n^{(k)}(a)}{k!}
={}&A_n(L)
 +\frac{nA_{n-1}(L)+n(n-1)A_{n-2}(L)}{2k}\notag\\
&+O_n\left(\frac{(1+L)^{n-1}}{k^2}\right)
 +O_{a_0,a_1,n}(k^2\rho^k).
\label{eq:fixed-index}
\end{align}
Terms with a negative Appell index are omitted.
To see the stated polynomial error, note that the analytic
remainder after the first correction vanishes at \(z=0\);
its coefficient convolution with \(e^{Lz}g(z)\) therefore has
degree at most \(n-1\) in \(L\). For \(n=0\) the first term is
exact and only the spectral error remains.
Higher powers of \(k^{-1}\) follow by the same coefficient procedure.
For example the \(k^{-2}\) numerator is
\[
\frac1{24}\left(
3\fall{n}{4}A_{n-4}
+2\fall{n}{3}A_{n-3}
-3\fall{n}{2}A_{n-2}
-2nA_{n-1}\right).
\]

At \(a=1\), \eqref{eq:stirling-product} identifies the first spectral
term with first-kind Stirling numbers. These are also the cycle-count
coefficients of permutations. The reciprocal-gamma factor is thus a
classical mechanism, closely related to mod-Poisson asymptotics
\cite{KowalskiNikeghbali}. The contribution to this programme is the
explicit spectral transfer, the uniform second coefficient estimate,
and the exact rational enclosure, not a claim that the
reciprocal-gamma phenomenon itself is new.

\section{Regularized modular rows and all derivative evaluations}
\label{sec:modular}

The modular draft obtains even-weight Eisenstein series by pairing
polygamma values along lattice rows. Its weight-two question requires
care: the paired rows converge rapidly, although the ungrouped lattice
sum is not absolutely convergent. We first prove an explicit bound for
every differentiated row. We then determine the change of summation
order directly, before using any modular transformation law. Finally,
Ramanujan's differential system evaluates all the resulting derivative
rows at the square and hexagonal CM points.

Throughout this section, let
\[
\tau\in\mathbb H:=\{z\in\C:\Impart z>0\},\qquad
q=e^{2\pi\ii\tau},\qquad r=|q|=e^{-2\pi\Impart\tau}<1,
\qquad \rho=\tau_{\mathrm{hex}}=\frac{1+\ii\sqrt3}{2}.
\]
For even \(w\ge2\) and an integer \(j\ge0\), put
\begin{equation}
\mathcal T_{w,j}(\tau)
=\sum_{n\ge1}n^j\left[
 \psi^{(w-1+j)}(n\tau)
 +(-1)^j\psi^{(w-1+j)}(1-n\tau)\right].
\label{eq:modular-T-definition}
\end{equation}
Convergence will follow from the next theorem. The sign on the reflected
term is part of the definition; differentiating that term changes its
sign at every step.

\subsection{Rational row formulas and explicit tail bounds}

The negative-index polylogarithms appearing below are elementary
rational functions. For \(d\ge0\), define
\[
R_d(t)=\Li_{-d}(t),\qquad
R_0(t)=\frac{t}{1-t},\qquad
R_{d+1}(t)=t\frac{d}{dt}R_d(t).
\]
These relations are identities of rational functions, since on
\(|t|<1\) they follow by differentiating
\(R_d(t)=\sum_{\ell\ge1}\ell^d t^\ell\). In particular,
\[
R_1(t)=\frac{t}{(1-t)^2},\qquad
R_2(t)=\frac{t(1+t)}{(1-t)^3},\qquad
R_3(t)=\frac{t(1+4t+t^2)}{(1-t)^4}.
\]

\begin{theorem}[Rational rows and an explicit remainder]
\label{thm:modular-rows-tail}
For \(w,j,\tau\) as above,
\begin{align}
&\psi^{(w-1+j)}(n\tau)
 +(-1)^j\psi^{(w-1+j)}(1-n\tau)
 =(2\pi\ii)^{w+j}\Li_{1-w-j}(q^n),
\label{eq:modular-single-row}\\
&\mathcal T_{w,j}(\tau)
 =(2\pi\ii)^{w+j}\sum_{n\ge1}n^j\Li_{1-w-j}(q^n).
\label{eq:modular-rational-rows}
\end{align}
The series converges normally on every compact subset of \(\mathbb H\).
If \(\mathcal T_{w,j}^{[N]}\) is its sum through \(n=N\), where
\(N\ge0\), then
\begin{equation}
\left|\mathcal T_{w,j}(\tau)-\mathcal T_{w,j}^{[N]}(\tau)\right|
\le (2\pi)^{w+j}(N+1)^j
 \frac{\Li_{-j}(r)}{r}\,
 \Li_{1-w-j}\!\left(r^{N+1}\right).
\label{eq:modular-tail-bound}
\end{equation}
Thus the bound contains only rational functions of \(r\) and
\(r^{N+1}\), together with an explicit power of \(\pi\).
\end{theorem}

\begin{proof}
For \(\Impart z>0\), splitting the absolutely convergent bilateral
power sum into its nonnegative and negative indices gives
\begin{equation}
\sum_{m\in\Z}(m+z)^{-w}
=\frac{\psi^{(w-1)}(z)+\psi^{(w-1)}(1-z)}{(w-1)!}.
\label{eq:modular-bilateral-polygamma}
\end{equation}
Here evenness of \(w\) is essential. The classical cotangent expansion
and its partial-fraction interpretation~\cite{DLMF} read
\[
\pi\cot(\pi z)=-\pi\ii
 -2\pi\ii\sum_{\ell\ge1}e^{2\pi\ii\ell z},\qquad
\pi\cot(\pi z)=\lim_{M\to\infty}\sum_{m=-M}^{M}\frac1{m+z}.
\]
Differentiate \(w-1\) times. On compact subsets of the upper
half-plane the exponential series and its derivatives converge
uniformly. After the first derivative, the bilateral series and all
its subsequent derivatives also converge normally. Since
\((-1)^{w-1}=-1\), the two minus signs cancel and yield
\[
\psi^{(w-1)}(z)+\psi^{(w-1)}(1-z)
=(2\pi\ii)^w\sum_{\ell\ge1}\ell^{w-1}e^{2\pi\ii\ell z}.
\]
Differentiating another \(j\) times with respect to \(z\) proves
\eqref{eq:modular-single-row}. Multiplication by \(n^j\) gives the
summand in \eqref{eq:modular-rational-rows}.

To prove the remainder bound, expand the rational polylogarithm in
its nonnegative majorizing series. The absolute value of the omitted
rows is at most
\begin{equation}
(2\pi)^{w+j}\sum_{\ell\ge1}\ell^{w+j-1}
                  \sum_{n>N}n^j r^{n\ell}.
\label{eq:modular-positive-majorant}
\end{equation}
Write \(n=N+1+a\). For \(a\ge0\) and \(\ell\ge1\),
\[
n^j\le(N+1)^j(a+1)^j,\qquad r^{a\ell}\le r^a.
\]
Consequently
\[
\sum_{n>N}n^j r^{n\ell}
\le (N+1)^j r^{(N+1)\ell}
       \sum_{a\ge0}(a+1)^j r^a
=(N+1)^j r^{(N+1)\ell}\frac{\Li_{-j}(r)}{r}.
\]
Substitution in \eqref{eq:modular-positive-majorant} proves
\eqref{eq:modular-tail-bound}. In particular, the majorant is finite.
On a compact subset of \(\mathbb H\), one has \(r\le r_0<1\), and
each of its power series with nonnegative coefficients is bounded by
its value at \(r_0\). This proves normal convergence, for every fixed
\(w,j\), and justifies all subsequent termwise differentiations.
\end{proof}

For \(j=0\), the tail formula simplifies to
\[
\left|\mathcal T_{w,0}-\mathcal T_{w,0}^{[N]}\right|
\le\frac{(2\pi)^w}{1-r}\Li_{1-w}\!\left(r^{N+1}\right).
\]
The case \(w=2\), which will be particularly useful, is
\begin{equation}
\left|\mathcal T_{2,0}-\mathcal T_{2,0}^{[N]}\right|
\le\frac{4\pi^2r^{N+1}}
 {(1-r)(1-r^{N+1})^2}.
\label{eq:modular-weight-two-tail}
\end{equation}
Theorem~\ref{thm:modular-rows-tail} also proves
\begin{equation}
\frac{d^j}{d\tau^j}\mathcal T_{w,0}(\tau)
=\mathcal T_{w,j}(\tau).
\label{eq:modular-termwise-derivatives}
\end{equation}
Computationally, \eqref{eq:modular-single-row} is preferable to
evaluating its two polygamma terms separately. Their algebraic
asymptotic contributions cancel, leaving the exponentially small
paired row. The rational formula performs that cancellation exactly
before numerical evaluation.

\subsection{The weight-two ordering correction from telescoping}

For \(w\ge4\), the ordinary absolutely convergent lattice series is
\[
G_w(\tau)=\sum_{(m,n)\in\Z^2\setminus\{(0,0)\}}
                  (m+n\tau)^{-w}
=2\zeta(w)+\frac{2}{(w-1)!}\mathcal T_{w,0}(\tau).
\]
For \(w=2\), specify the order explicitly:
\begin{align}
G_2(\tau)&=\sum_{n\in\Z}
        \left(\sum_{m\in\Z}'(m+n\tau)^{-2}\right)
 =\frac{\pi^2}{3}+2\mathcal T_{2,0}(\tau),
\label{eq:modular-G2-definition}\\
C_2(\tau)&=\sum_{m\in\Z}
        \left(\sum_{n\in\Z}'(m+n\tau)^{-2}\right).
\label{eq:modular-C2-definition}
\end{align}
The prime omits only \((m,n)=(0,0)\). In the first expression, each
inner row and the outer series of completed rows converge absolutely.
The reverse expression exists as well: changing the names and signs
of the one-dimensional indices gives
\begin{equation}
C_2(\tau)=\tau^{-2}G_2(-1/\tau).
\label{eq:modular-C2-coordinate}
\end{equation}
This is a coordinate identity for the prescribed iterated sums, not
an assumption of modularity. The series of absolute values over all
lattice points diverges, so an interchange of the two orders still
requires a calculation.

\begin{theorem}[Summation-order anomaly]
\label{thm:modular-ordering}
For every \(\tau\in\mathbb H\),
\begin{equation}
C_2(\tau)=G_2(\tau)-\frac{2\pi\ii}{\tau}.
\label{eq:modular-anomaly}
\end{equation}
Consequently
\begin{equation}
G_2(-1/\tau)=\tau^2G_2(\tau)-2\pi\ii\tau,
\qquad G_2(\tau+1)=G_2(\tau).
\label{eq:modular-G2-transform}
\end{equation}
\end{theorem}

\begin{proof}
We prove the difference of the two orders by replacing their common
part with an absolutely convergent series. Put
\[
h(z)=\frac1z-\frac1{z+1},\qquad
f(z)=\frac1{z^2}-h(z)=\frac1{z^2(z+1)},
\]
and define
\[
F(\tau)=\sum_{\substack{n\in\Z\setminus\{0\}\\m\in\Z}}
                    f(m+n\tau).
\]
No pole occurs in this sum. The real-linear map
\((m,n)\mapsto m+n\tau\) is invertible, so for some \(c_\tau>0\),
\(|m+n\tau|\ge c_\tau\sqrt{m^2+n^2}\).
Outside a finite set, \(|f(z)|\le2|z|^{-3}\); hence \(F\) converges
absolutely. It can therefore be summed in either order.

For each fixed \(n\ne0\), the series \(\sum_m h(m+n\tau)\) is
absolutely convergent and its symmetric partial sum telescopes:
\[
\sum_{m=-M}^{M}h(m+n\tau)
=\frac1{-M+n\tau}-\frac1{M+1+n\tau}\longrightarrow0.
\]
Separating the row \(n=0\) proves
\begin{equation}
G_2(\tau)=\frac{\pi^2}{3}+F(\tau).
\label{eq:modular-G2-absolute-part}
\end{equation}

For the reverse order, define for integer \(m\)
\[
A(m)=\lim_{K\to\infty}
              \sum_{0<|n|\le K}\frac1{m+n\tau}.
\]
This symmetric limit exists: pairing \(n\) and \(-n\) gives an
\(O(n^{-2})\) summand. The cotangent partial-fraction formula gives
\begin{equation}
A(m)=\frac{\pi}{\tau}\cot\!\left(\frac{\pi m}{\tau}\right)-\frac1m
\quad(m\ne0),\qquad A(0)=0.
\label{eq:modular-A-cotangent}
\end{equation}
The value at zero is also the removable limit of the expression on
the right. For each fixed \(m\), the \(n\)-series of \(h(m+n\tau)\)
is absolutely convergent. Taking its symmetric partial sums shows
\[
\sum_{n\ne0}h(m+n\tau)=A(m)-A(m+1).
\]
It follows that
\begin{align*}
C_2(\tau)-\frac{\pi^2}{3}-F(\tau)
&=\lim_{M\to\infty}\sum_{m=-M}^{M}\sum_{n\ne0}h(m+n\tau)\\
&=\lim_{M\to\infty}\bigl(A(-M)-A(M+1)\bigr).
\end{align*}
The extraction of \(F\) here is legitimate by its absolute
convergence; no rearrangement of the original double power sum has
been made.

Since \(\Impart(1/\tau)<0\),
\[
\cot(\pi m/\tau)\longrightarrow\ii\quad(m\to+\infty),
\qquad
\cot(\pi m/\tau)\longrightarrow-\ii\quad(m\to-\infty).
\]
Equation~\eqref{eq:modular-A-cotangent} therefore gives
\(A(-M)-A(M+1)\to-2\pi\ii/\tau\). Combining this with
\eqref{eq:modular-G2-absolute-part} proves
\eqref{eq:modular-anomaly}. Equation~\eqref{eq:modular-C2-coordinate}
now gives the first transformation in
\eqref{eq:modular-G2-transform}. Periodicity follows directly from
the \(q\)-series in \eqref{eq:modular-rational-rows}.
\end{proof}

The anomaly is classical; Romik and Scherer discuss it together with
more general shape-dependent summation procedures~\cite{RomikScherer}.
The proof above fixes both its sign and the normalization needed for
the draft's polygamma sums.

\begin{corollary}[The missing weight-two CM rows]
\label{cor:modular-weight-two-CM}
One has
\begin{align}
G_2(\ii)&=\pi,&G_2(\rho)&=\frac{2\pi}{\sqrt3},
\label{eq:modular-G2-CM}\\
\mathcal T_{2,0}(\ii)&=\frac\pi2-\frac{\pi^2}{6},&
\mathcal T_{2,0}(\rho)&=\frac\pi{\sqrt3}-\frac{\pi^2}{6}.
\label{eq:modular-trigamma-CM}
\end{align}
\end{corollary}
\begin{proof}
At \(\tau=\ii\), the first transformation in
\eqref{eq:modular-G2-transform} gives
\(G_2(\ii)=-G_2(\ii)+2\pi\).
For \(\rho\), use \(-1/\rho=\rho-1\) and periodicity. The same
transformation becomes
\[
(\rho^2-1)G_2(\rho)=2\pi\ii\rho.
\]
Since \(\rho^2-1=\ii\sqrt3\rho\), this proves the second value in
\eqref{eq:modular-G2-CM}. Substitution into
\eqref{eq:modular-G2-definition} proves
\eqref{eq:modular-trigamma-CM}.
\end{proof}

\subsection{CM seeds and the complete differential tower}

Normalize the Eisenstein series by
\[
E_2=\frac{3G_2}{\pi^2},\qquad
E_4=\frac{G_4}{2\zeta(4)},\qquad
E_6=\frac{G_6}{2\zeta(6)}.
\]
The seed values, in precisely this normalization, are
\begin{align}
(E_2(\ii),E_4(\ii),E_6(\ii))
&=\left(\frac3\pi,
 \frac{3\Gamma(1/4)^8}{64\pi^6},0\right),
\label{eq:modular-square-seeds}\\
(E_2(\rho),E_4(\rho),E_6(\rho))
&=\left(\frac{2\sqrt3}\pi,0,
 \frac{27\Gamma(1/3)^{18}}{512\pi^{12}}\right).
\label{eq:modular-hex-seeds}
\end{align}
We have already proved the weight-two entries. The zero entries follow
by rotating the absolutely convergent lattices. Multiplication by
\(\ii\) fixes the square lattice and sends its weight-six sum to its
negative; multiplication by \(\rho\) fixes the hexagonal lattice and
multiplies its weight-four sum by \(\rho^{-4}\ne1\).

For completeness, the classical period integrals fix the nonzero
entries without numerical identification. The standard inverse
Weierstrass integral~\cite{DLMF}, with
\(g_2=60G_4\) and \(g_3=140G_6\), gives positive real periods
\begin{align*}
p_{\mathrm{sq}}
&=2\int_1^\infty\frac{dx}{\sqrt{4x^3-4x}}
 =2\int_0^1\frac{du}{\sqrt{1-u^4}}
 =\frac12 B\!\left(\frac14,\frac12\right)
 =\frac{\Gamma(1/4)^2}{2\sqrt{2\pi}},\\
p_{\mathrm{hex}}
&=2\int_1^\infty\frac{dx}{\sqrt{4x^3-4}}
 =\frac13 B\!\left(\frac16,\frac12\right)
 =\frac{\Gamma(1/3)^3}{2^{4/3}\pi}.
\end{align*}
The corresponding lattices are
\(p_{\mathrm{sq}}(\Z+\ii\Z)\) for \((g_2,g_3)=(4,0)\), and
\(p_{\mathrm{hex}}(\Z+\rho\Z)\) for \((g_2,g_3)=(0,4)\).
Their square and equianharmonic shapes follow from the rotations of
the respective cubics; the real period integral fixes the scale and
orientation. The first substitution above is \(x=u^{-2}\); in the
second beta integral use \(t=x^{-3}\). Euler's reflection and
duplication formulas give the displayed gamma evaluations.
Scaling the lattice invariants consequently yields
\[
G_4(\ii)=\frac{p_{\mathrm{sq}}^4}{15}
 =\frac{\Gamma(1/4)^8}{960\pi^2},\qquad
G_6(\rho)=\frac{p_{\mathrm{hex}}^6}{35}
 =\frac{\Gamma(1/3)^{18}}{8960\pi^6}.
\]
Using \(2\zeta(4)=\pi^4/45\) and
\(2\zeta(6)=2\pi^6/945\) proves the remaining entries of
\eqref{eq:modular-square-seeds}--\eqref{eq:modular-hex-seeds}.

Ramanujan's classical differential identities~\cite{Ramanujan} are
\begin{equation}
\theta E_2=\frac{E_2^2-E_4}{12},\qquad
\theta E_4=\frac{E_2E_4-E_6}{3},\qquad
\theta E_6=\frac{E_2E_6-E_4^2}{2},\qquad
\theta=\frac1{2\pi\ii}\frac{d}{d\tau}.
\label{eq:modular-Ramanujan-system}
\end{equation}
Encode these identities by the rational derivation
\begin{equation}
\mathcal V
=\frac{X^2-Y}{12}\frac{\partial}{\partial X}
 +\frac{XY-Z}{3}\frac{\partial}{\partial Y}
 +\frac{XZ-Y^2}{2}\frac{\partial}{\partial Z},
\qquad P_0=X,\quad P_{j+1}=\mathcal V P_j.
\label{eq:modular-P-recurrence}
\end{equation}
For example,
\begin{align*}
P_1&=\frac{X^2-Y}{12},\\
P_2&=\frac{X^3-3XY+2Z}{72},\\
P_3&=\frac{X^4-6X^2Y+8XZ-3Y^2}{288}.
\end{align*}

\begin{theorem}[All differentiated weight-two rows]
\label{thm:modular-all-jets}
Set \(\mathcal S_j(\tau)=\mathcal T_{2,j}(\tau)\). For every
\(\tau\in\mathbb H\),
\begin{align}
\mathcal S_0(\tau)&=\frac{\pi^2}{6}\bigl(E_2(\tau)-1\bigr),
\label{eq:modular-S0}\\
\mathcal S_j(\tau)&=\frac{\pi^2}{6}(2\pi\ii)^j
 P_j(E_2(\tau),E_4(\tau),E_6(\tau)),\qquad j\ge1.
\label{eq:modular-all-jets}
\end{align}
In particular, substitution of
\eqref{eq:modular-square-seeds}--\eqref{eq:modular-hex-seeds}
evaluates every \(\mathcal S_j(\ii)\) and
\(\mathcal S_j(\rho)\) by a finite polynomial calculation over
\(\Q\) and the indicated gamma periods.
\end{theorem}

\begin{proof}
Equation~\eqref{eq:modular-G2-definition} and the definition of
\(E_2\) give \eqref{eq:modular-S0}. By the chain rule and
\eqref{eq:modular-Ramanujan-system}, every polynomial
\(P\in\Q[X,Y,Z]\) satisfies
\[
\theta\bigl(P(E_2,E_4,E_6)\bigr)
=(\mathcal V P)(E_2,E_4,E_6).
\]
Induction from \(P_0=X\) therefore gives
\(\theta^jE_2=P_j(E_2,E_4,E_6)\).
Differentiate \eqref{eq:modular-S0} \(j\ge1\) times and use
\eqref{eq:modular-termwise-derivatives}. The constant term disappears,
and \(d/d\tau=2\pi\ii\theta\) contributes the factor
\((2\pi\ii)^j\). This proves \eqref{eq:modular-all-jets}; the CM
specializations use the established seed values.
\end{proof}

For clarity, the first two differentiated evaluations are
\begin{align}
\sum_{n\ge1}n\left[\psi''(n\ii)-\psi''(1-n\ii)\right]
 &=\frac{\ii\pi}{4}
   -\frac{\ii\Gamma(1/4)^8}{768\pi^3},
\label{eq:modular-S1-square}\\
\sum_{n\ge1}n\left[\psi''(n\rho)-\psi''(1-n\rho)\right]
 &=\frac{\ii\pi}{3},
\label{eq:modular-S1-hex}\\
\sum_{n\ge1}n^2\left[\psi'''(n\ii)+\psi'''(1-n\ii)\right]
 &=-\frac\pi4+\frac{\Gamma(1/4)^8}{256\pi^3},
\label{eq:modular-S2-square}\\
\sum_{n\ge1}n^2\left[\psi'''(n\rho)+\psi'''(1-n\rho)\right]
 &=-\frac{2\sqrt3\pi}{9}
   -\frac{\Gamma(1/3)^{18}}{1024\pi^8}.
\label{eq:modular-S2-hex}
\end{align}
These follow by substituting the seed values into \(P_1,P_2\), so
their derivation contains no integer-relation search.

\begin{corollary}[Polynomial degree bounds for the CM row formulas]
\label{cor:modular-CM-degree}
Put
\(A=\Gamma(1/4)^8/\pi^4\) and
\(B=\Gamma(1/3)^{18}/\pi^9\). For each \(j\ge1\), there are
polynomials \(U_j\in\Q[t]\) and
\(V_j\in\Q(\sqrt3)[t]\) such that
\[
\frac{\mathcal S_j(\ii)}{\ii^j\pi}=U_j(A),\qquad
\frac{\mathcal S_j(\rho)}{\ii^j\pi}=V_j(B),
\]
with
\[
\deg U_j\le\left\lfloor\frac{j+1}{2}\right\rfloor,
\qquad
\deg V_j\le\left\lfloor\frac{j+1}{3}\right\rfloor.
\]
\end{corollary}
\begin{proof}
Give \(X,Y,Z\) weights \(2,4,6\). Each term of
\(\mathcal V\) raises weight by two, so \(P_j\) is homogeneous
of weight \(2j+2\). At the square point, the surviving monomials
\(X^aY^b\) have \(a+2b=j+1\). On inserting the seeds into
\eqref{eq:modular-all-jets} and dividing by \(\ii^j\pi\), each
becomes a rational multiple of \(A^b\). At the hexagonal point,
the surviving monomials \(X^aZ^c\) satisfy \(a+3c=j+1\), and
the same substitution produces an element of
\(\Q(\sqrt3)B^c\). The stated degree bounds follow.
\end{proof}

\subsection{What this resolves and how it is verified}

Corollary~\ref{cor:modular-weight-two-CM} resolves the weight-two
row-sum question posed in the modular draft. Theorem~\ref{thm:modular-all-jets}
extends it to all polynomially weighted derivative rows, with the
explicit error estimate of Theorem~\ref{thm:modular-rows-tail}.
These are evaluations of specified infinite aggregates. They do not
evaluate an isolated component such as \(\Impart\psi'(\ii)\), nor
do they prove that such a component is an Eichler integral. Any proposed
identification of that isolated constant needs its own exact transform
or functional equation.

The accompanying \path{code/modular_rows.py} constructs the rational
functions \(R_d\) and the polynomials \(P_j\) by exact arithmetic.
It checks individual rows against a separate complex-polygamma
implementation, compares the CM formulas for \(j=0,\ldots,4\) with
their rapidly convergent row sums, and checks representative tails.
The output \path{data/modular_rows.json} records the exact coefficients,
analytic truncation bounds, and numerical residuals. Its decimal
calculations use ordinary high-precision arithmetic; they are numerical
checks, not outward-rounded interval certificates. A machine enclosure
can be obtained by evaluating the rational formulas and
\eqref{eq:modular-tail-bound} with directed interval arithmetic. The
identities themselves are proved above and do not depend on the
numerical checks.

\section{Reproducibility, certification, and integration}
\label{sec:reproducibility}

\subsection{What the computations establish}

The proofs above establish the theorems; the computations independently
check normalization, signs, low-order specializations, and implementation.
They are deliberately recorded in three distinct categories.
\begin{enumerate}
\item \textbf{Exact algebra.} Partial-fraction identities are checked
as rational polynomial identities. Distribution and reflection matrices
are reduced over exact rational arithmetic. The modular differential
polynomials retain rational coefficients.
\item \textbf{Analytically certified enclosures.} The intervals in
Corollary~\ref{cor:rational} have exact rational endpoints. Their validity
comes from the proved spectral-tail bound, not a comparison of decimal
approximations.
\item \textbf{Independent high-precision checks.} Complex quadrature,
Hurwitz-zeta evaluation, and polygamma evaluation test the formulas by
different representations. These use ordinary arbitrary-precision
floating-point arithmetic, without outward rounding. They provide
strong implementation checks but are not interval certificates.
\end{enumerate}

The shipped default runs have the following scope.
\begin{center}
\small
\begin{tabular}{>{\raggedright\arraybackslash}p{0.27\textwidth}p{0.64\textwidth}}
\toprule
Artifact & Recorded verification\\
\midrule
Mixed doubles and Euler sums &
64 exact partial-fraction checks; 18 mixed integral comparisons;
the four explicit reductions; five Euler values; six generator checks;
three pole-residue checks. Working precision: 80 decimal digits.\\
\addlinespace
Distributions and bridges &
1,734 exact matrix checks through \(q=60\), including anchored and
reflected quotients and all-divisor/prime-row agreement.
72 analytic checks of normalized bridges through derivative order five,
at 65 decimal digits.\\
\addlinespace
Parameter asymptotics &
78 rational enclosures, including 72 strictly positive sign certificates;
72 independent exact coefficient comparisons; 36 numerical
spectral-tail checks at 75 decimal digits; the profile figure.\\
\addlinespace
Modular rows &
Four independent single-row comparisons, ten CM jet comparisons through
order four, and four explicit tail checks at 90 decimal digits;
exact differential polynomials through order five.\\
\bottomrule
\end{tabular}
\end{center}

The maximum equality residual in the mixed/Euler run was
approximately \(1.71\times10^{-80}\), against the declared tolerance
\(10^{-60}\). Pole residues use their own finite-difference error
criterion and are not included in that equality statistic.
The source manifest verifies all 40 retrieved text files against their
Git blob identifiers and records additional SHA-256 hashes.
This is provenance for the reviewed snapshot, not a claim that every
statement in all 31 reports has been checked.

\subsection{A certificate format that survives integration}

The exact objects should remain separate from their numerical
evaluations. For a linear relation, store the convention and ordered
symbol list, the integer or rational coefficient vector, the source
law identifiers, and the exact combination of law rows that produces
the vector. An integer-relation search may propose that vector, but
does not supply its law certificate.

For a finite-dimensional quotient, store the relation system and an
exact normal-form map in addition to the rank. A rank number without
the rows and their interpretation is not enough to reconstruct the
claim. The character basis in Section~\ref{sec:distribution} provides
a particularly compact symbolic normal form across infinitely many
denominators. Its general proof is stronger than a larger finite table.

For numerical values, store the representation, parameters, precision,
truncation, and the actual error statement. Where a proved analytic
majorant is evaluated in ordinary floating point, label the result
accordingly; an outward-rounded evaluation is an additional step.
For the rational intervals, retain numerator and denominator rather
than only a decimal rendering.

The package is self-contained at the source level: the master
\src{article.tex} includes the section files, the bibliography is
included, the figure source and data are present, and a build command
is supplied in \src{README.txt}. The Python scripts use
\src{mpmath}, \src{sympy}, \src{numpy}, and \src{matplotlib}.
No absent upstream value store is needed to reproduce these results.
The proposed corrections are supplied as a separate register so that
maintainers can review and apply them to the original articles.
The original corpus is neither duplicated nor silently modified.

\section{Further research: precise questions and proposed routes}
\label{sec:questions}

The corrections change which problems are worth pursuing. In
particular, the inverse-argument mixed family no longer supplies the
proposed new directions. A productive continuation should use exact
relations to remove those directions before allocating numerical
effort to the genuinely unresolved quotient.

\subsection{Immediate structural projects}

\paragraph{1. A convention-safe cyclotomic reduction engine.}
Can the partial-fraction map of Theorem~\ref{thm:mixed} be combined
with shuffle, stuffle, distribution, conjugation, and the proven
parity theorem into a single exact reducer at fixed weight and level?
Start with levels three, four, and six, retain an explicit ordered
alphabet, and distinguish products from the linear depth filtration.
Every emitted relation should carry a rational certificate.
The first target is a reproducible spanning set and formal corank
through a stated weight, not an arithmetic dimension theorem.
The four false survivors make useful regression examples.

\paragraph{2. The remaining parity of alternating harmonic sums.}
For
\[
S_{2m}=\sum_{n\ge0}\frac{(-1)^nH_n}{(2n+1)^{2m}},
\]
what is the smallest natural motivic or formal level-four space that
contains the family? The beta-integral argument evaluates the other
parity but does not prove nonmembership here. A concrete next step is
to compute a depth-graded coaction or an exact formal quotient for the
first few weights, identify the image of \(S_{2m}\), and distinguish
the resulting motivic statement from any period-independence
assumption. Xu and Wang's wider Euler-sum framework is a natural
starting point~\cite{XuWang}.

\paragraph{3. Higher-depth inverse-color reductions.}
The substitution \(n=m+h\) works especially well because the inverse
arguments leave a single phase \(z^h\). For triples such as
\(\Li_{a,b,c}(z,z^{-1},1)\), which phase patterns lead to a similar
gap-variable reduction, and which genuinely retain two interacting
gaps? A first target is an exact partial-fraction identity at depth
three with a complete boundary-convergence proof. Merely increasing
the depth-two search basis cannot establish that a triple lies in it.

\paragraph{4. Integral weighted distributions.}
The character proof works over a \(\C\)-algebra, where averaging over
\((\Z/q\Z)^\times\) and diagonalizing characters are available.
What is the corresponding integral module over
\(\Z[t_p:p\mid q]\), and which torsion appears after reflection,
anchoring, or specialization? The distinction matters if a future
certificate system is to use integral normal forms without inverting
group orders. Smith normal forms for prime powers and two-prime
denominators would supply useful conjectures; an induction or a
distribution resolution is needed for the general theorem.
The connection to universal distributions should be retained
\cite{Kubert,Ouyang}.

\subsection{Asymptotic and analytic projects}

\paragraph{5. Beyond the logarithmic index range.}
Theorem~\ref{thm:profile} is uniform for \(n\le C\log(k/a)\).
What happens when \(n/\log k\to\infty\)? The exact coefficient
representation suggests a saddle-point equation for
\[
e^{-z\log a}\prod_{j=1}^{k}(1+z/j),
\]
but the appropriate saddle may move far enough that fixed-disk
gamma-ratio estimates no longer apply. The competing \(m\ge1\)
spectral terms also need uniform control.
A concrete first extension is a proved range
\(n=o((\log k)^2)\), with an error estimate that explicitly describes
where the two-term profile ceases to be useful. That range is a
research target here, not an established theorem.

\paragraph{6. Moving Hurwitz parameter and spectral competition.}
The present compact-\(a\) hypothesis gives
\((a/(a+1))^k\) as an exponentially small ratio.
For \(a=a(k)\) comparable with \(k\), this ratio need not tend to zero.
Can one derive a profile involving several spectral terms, or an
integral limit of them? The separate regimes
\(a=o(k)\), \(a\asymp k\), and \(a\gg k\) should be analyzed using
the exact expansion before proposing a single universal formula.
This would connect the parameter tower with large-parameter
Hurwitz-zeta asymptotics.

\paragraph{7. Fully effective asymptotic constants and optimal bounds.}
Lemma~\ref{lem:coefficient} has explicit constants, while
Theorem~\ref{thm:profile} packages the fixed-disk bounds into \(O\)-terms.
An effective version could expose a computable \(k_0(a_0,a_1,C)\)
for the sign theorem. A separate optimization problem is to minimize
\eqref{eq:spectral-tail} over \(R\) and the truncation \(M\), or to
compare the integer-radius rational certificate with a rational
upper approximation to the optimal real radius.
The success criterion is a rigorously narrower interval at a known
computational cost, not simply more printed digits.

\paragraph{8. Resonant jet coordinates and numerical conditioning.}
The full distribution module is free even where top-denominator
coordinates vanish. How should a symbolic or numerical algorithm
choose conductor-level coordinates so that the conversion remains
well-conditioned near a resonance? The vanishing-order formula in
Section~\ref{sec:distribution} provides a local answer, including the
number of lost jet coefficients. A useful next deliverable is a
canonical basis-selection algorithm with explicit transition matrices,
followed by a study of conditioning at large denominators.
No extra numerical constants should be inferred from a coordinate
singularity.

\paragraph{9. Global continuation of normalized derivative bridges.}
Section~\ref{sec:bridges} gives a local germ after removing a simple
trivial zero. How far can each germ be continued before an actual
zero of the \(L\)-function or a known elementary factor becomes an
obstruction? One can formulate zero-sensitive domains for the
logarithmic jets, derive Bell-polynomial formulas for raw derivatives,
and build exact symbolic identities for higher orders.
For imprimitive characters, Euler-factor zeros should be removed
explicitly and compared with the distribution resonances.
The root-number cancellation in the local identity should not be
misread as a global zero-free theorem.

\subsection{Arithmetic and modular questions}

\paragraph{10. The unresolved arithmetic content of gamma quotients.}
The formal gamma quotient has a proved rank, and every law-derived
identity has an exact certificate. Which additional arithmetic
relations could remain beyond that system? This is the point at
which Rohrlich-type conjectures enter. A useful project would separate
three outputs: unconditional law certificates, conditional
completeness statements with the exact conjecture written out, and
finite bounded-height exclusions using certified numerical inputs.
The Koblitz--Ogus sufficient criterion must not be promoted to the
missing converse; see the careful distinction in
Anderson--Brownawell--Papanikolas~\cite{ABP}.

\paragraph{11. Isolated imaginary polygamma values at CM points.}
The modular row theorems evaluate weighted aggregates. They do not
individually evaluate \(\Impart\psi^{(1)}(\ii)\).
Does this isolated quantity belong to a natural Eichler-integral,
iterated-integral, or regulator framework with a provable modular
transformation law? The immediate target is the transformation law
itself, including its nonholomorphic or boundary terms if necessary.
A Lambert-series resemblance alone is not sufficient evidence.
Once such a law is established, it may suggest period classes and
meaningful relations among different CM arguments.

\paragraph{12. Beyond the square and hexagonal lattices.}
At other imaginary-quadratic parameters, the same Ramanujan
recurrence computes all row derivatives from \(E_2,E_4,E_6\).
Can the algebraic CM data and periods be certified exactly, with
minimal polynomials, isolating regions, and compatible period
normalizations? A useful first target is several class-number-two
discriminants. A polynomial relation supplies an upper bound on
degree; exact degree additionally requires irreducibility and
identification of the intended root. This would replace the draft's
PSLQ degree guesses by checkable algebraic evidence.

\subsection{A sensible order of attack}

The first, fourth, seventh, and eighth projects have the clearest
near-term certification goals. They strengthen the mathematical
infrastructure without depending on a period conjecture. The second,
third, and ninth concern deeper structural reductions, where formal
or motivic calculations can establish substantial results before
arithmetic independence is addressed. The remaining asymptotic and CM
projects require new analytic input.

The central methodological requirement is to keep three maps visible:
from exact functional identities to a formal quotient, from that
quotient to actual special values, and from those values to numerical
approximations. A proof about one map should not silently become a
claim about another.

\appendix
\section{Source-specific corrections and proof-status revisions}
\label{app:audit}

This appendix gives the principal proposed changes. The accompanying
\src{CORRECTIONS.txt} contains a fuller register, including exact
source line ranges and smaller notation repairs. All locators refer
to the commit on the title page. A false identity, an unproved claim,
and a correct statement with an incomplete proof are different
findings; they are identified separately below.

\subsection{The inverse-color mixed report}

\noindent\textbf{Source:}
\path{reports/eisenstein-gaussian-mixed-cuberoot-doubles__abd2e7116aa9.tex}.

The abstract, the table at source lines 170--184, and especially
\src{thm:parallel} at lines 219--226 claim two new constants in
each mixed tower. This is an \textbf{established error} relative to
the stated base family. Corollary~\ref{cor:four} gives explicit
single-data reductions of all four named candidates, and
Theorem~\ref{thm:mixed} eliminates any additional inverse-color
directions at every positive integral weight.

Replace the theorem by the all-weight quotient statement in
Section~\ref{sec:mixed}; update the abstract, table interpretation,
and any later dimension counts that use those directions.
The report's general numerical-evidence caveat does not repair a
false displayed theorem.
The claimed infinite asymptotic tail at lines 104--109 also needs a
finite truncation with a proved remainder, and the requested-digit
guarantee for Wynn acceleration at lines 111--118 needs an actual
error theorem. These are separate from the exact mixed reduction.
The missing upstream evaluator prevents a specific diagnosis of the
original search failure.

\subsection{Foundational multiple-polylogarithm conventions}

\noindent\textbf{Source:}
\path{articles/multiple-polylogarithms-introduction.tex}.

\paragraph{Convergence (lines 231--235).}
The stated absolute-convergence theorem is false already for
\(\Li_1(\ii)=\sum_{n\ge1}\ii^n/n\): the original series converges,
but its series of absolute values diverges.
A safe replacement in the stated positive-index setting is:
for \(|z_j|\le1\) and \(\Repart s_j\ge1\), absolute convergence
holds if \(\Repart s_1>1\) or \(|z_1|<1\).
More generally, strict inequalities
\(|z_1\cdots z_j|<1\) for every prefix give an interior domain of
absolute convergence. Boundary convergence and permissible
rearrangements should be stated separately. For positive integral
multiple-zeta indices, ordinary convergence is equivalent to
\(s_1\ge2\).

\paragraph{Iterated-integral orientation (lines 278--287).}
Use a consistent outermost-first recursion
\[
G(a_1,\ldots,a_w;z)=\int_0^z\frac{dt}{t-a_1}
                         G(a_2,\ldots,a_w;t),\qquad G(;z)=1.
\]
Its simplex is \(0<t_w<\cdots<t_1<1\), with \(a_j\) attached to
\(t_j\), and differentiation strips the \emph{first} letter.
The corresponding dictionary is
\begin{align*}
\Li_{s_1,\ldots,s_d}(z_1,\ldots,z_d)
=(-1)^dG(&0^{s_1-1},z_1^{-1},
          0^{s_2-1},(z_1z_2)^{-1},\ldots,\\
         &0^{s_d-1},(z_1\cdots z_d)^{-1};1).
\end{align*}
The arguments are inverse prefix products. The printed suffix
products are incompatible with this convention.
The same simplex reversal is needed in
\path{articles/gaussian-eisenstein-double-polylogs.tex},
\src{eq:G-def}--\src{eq:goncharov-dbl}, lines 141--146:
the printed orientation would make its proposed \(G(0,z^{-1};1)\)
representation diverge at the inner integration endpoint.

\paragraph{HPL recursion and signs (lines 491--555).}
A leading zero does not generally cause divergence:
\(H(0,1;z)=\Li_2(z)\) is a direct counterexample.
Use the recursion whenever its tail makes the integral convergent,
and define all-zero words separately by \(\log^w(z)/w!\).
With the source's convention
\(\zeta(\overline{s})=\sum_{n\ge1}(-1)^n/n^s\),
\[
H(0,-1;1)=-\Li_2(-1)=-\zeta(\overline2)=\frac{\pi^2}{12}.
\]
The sign in its identification with the barred zeta value must change.

\paragraph{Alphabet closure (lines 559--579).}
The full alphabet \(\{0,1,-1,\infty\}\) is not preserved by
every listed substitution. The weight-one counterexamples
\[
H(-1;1-z)=\log(2-z),\qquad
H(-1;z^2)=\log(1+z^2)
\]
introduce the singularities \(2\) and \(\pm\ii\), respectively.
The transformations \(-z\), \(1/z\), and
\((1-z)/(1+z)\) do preserve the singular set, with branches tracked.
Reflection \(1-z\) preserves the smaller \(\{0,1\}\) alphabet.
An algorithm acting on the full class must allow the necessary
enlarged alphabet.

\paragraph{Shuffle versus stuffle (lines 759--763).}
The correct identities are
\[
\Li_1(z)^2=2\Li_{1,1}(z,1)
          =2\Li_{1,1}(z,z)+\Li_2(z^2).
\]
The displayed formula with \(2\Li_{1,1}(z,z)\) alone is false:
that double starts at order \(z^3\), while the square starts at
order \(z^2\). The neighboring word example for
\(\zeta(2)\zeta(3)\) should shuffle the words \(01\) and \(001\),
of total weight five, rather than the printed weight-eight pair.

\paragraph{Unsupported numerical dimensions.}
The weight-five multiple-zeta table, lines 471--478, supplies
spanning reductions but not the rational independence of its two
generators. Replace the actual dimension assertion by an upper bound
or a stated conjectural value. Related sign and kernel repairs are
listed in the correction register; they do not invalidate every
low-weight identity in the introduction.

\subsection{Depth, parity, and proposed arithmetic bases}

\noindent\textbf{Sources:}
\path{articles/gaussian-multiple-polylog-depth.tex} and
\path{articles/gaussian-eisenstein-double-polylogs.tex}.

The definition of \(S_p\) in the first article, lines 128--132,
must require \(p\ge1\) for the ordinary series. At \(p=0\), its
terms fail to tend to zero.
Replace the \emph{precisely when} content of \src{thm:alternation},
lines 185--192, by the proved positive inclusion for every
\(S_{2m+1}\). Section~\ref{sec:euler} supplies its exact formula.
The source's \(S_5\) and \(S_7\) identities survive this change.
Its even-index nonmembership assertion is unproved; this was
already recognized in the intake README.

The uncolored analogy at line 351 includes
\(\zeta(1,2,1)\) and \(\zeta(1,1,2)\), which diverge as ordinary
positive-term multiple-zeta series. A regularized statement needs
an explicit regularization and a derivation.
Statements describing even beta values as new transcendental numbers
or a Champernowne canary as jointly algebraically independent are
not established by the supplied arguments.

In the second article, the dimension interpretations around
\src{thm:gauss-5} and \src{thm:eis-1} are \textbf{unsupported
arithmetic lower bounds}, not identities disproved here.
Replace them by the demonstrated spanning claims, or by an exact
formal corank if the complete matrix and certificate are supplied.
Products of depth-one and depth-two values lie in \(D^3\) in general:
\[
D^aD^b\subseteq D^{a+b}.
\]
The algebra generated by all depth-at-most-two values is a different
object from the linear filtration step \(D^2\).

The binomial depth-dimension formula near \src{thm:deligne},
lines 345--352, does not follow from free-Lie generation alone.
Word length in chosen Lie generators is not automatically MPL depth.
Even a positive ambient motivic dimension does not identify the
images of the three selected periods as nonzero or independent.
Adding a period conjecture cannot supply these missing motivic
calculations. Retain the proved reductions and reopen the specific
depth-three classification.

\subsection{Gamma laws, formal ranks, and complexity}

\noindent\textbf{Source:}
\path{articles/gamma-lattice-and-certificates.tex}.

The arithmetic completeness attributed to Koblitz--Ogus in
\src{law:ko}, lines 133--137, is not an available theorem of that
scope. The sufficient algebraicity criterion and the conjectural
converse must be distinguished; the intake README already flags
this issue. A primary discussion of the distinction is
\cite[\S1.4.2]{ABP}.

The unconditional replacement is:
\emph{The exact reducer is sound and complete for membership in the
rational span of the encoded reflection and multiplication rows.
Its free columns form a basis of that formal quotient.
Completeness for all multiplicative relations among actual gamma
values is not asserted.}
Section~\ref{sec:distribution} gives the formal anchored reflected
rank \(\varphi(N)/2\) for \(N\ge3\); at \(N=120\), this is 16.
The explicit identities accompanied by valid law certificates remain
valid.

The raw all-divisor row construction, lines 105--121, emits
\[
\left\lfloor\frac{N-1}{2}\right\rfloor
 +\sum_{\substack{d\mid N\\d\ge2}}\frac Nd
=\left\lfloor\frac{N-1}{2}\right\rfloor+\sigma_1(N)-N.
\]
This is not uniformly \(O(N)\), since \(\sigma_1(N)/N\) is
unbounded. State the exact count or a valid general upper bound;
row dependency does not retroactively change the raw output count.

\subsection{Stieltjes ranks and the complex-character correction}

\noindent\textbf{Sources:}
\path{articles/stieltjes-parameter-derivative-tower.tex},
\path{articles/stieltjes-antiderivative-ladder.tex}, and
\path{reports/stieltjes-derivative-relations__41afbe98f668.md}.

The derivative article's \src{thm:rank}, lines 195--212, states
an all-\(q,k\) conclusion with finite evidence.
The antiderivative article's \src{prop:jetrank}, lines 946--975,
more carefully states its finite range and conjectures the extension.
The missing all-denominator proof is now supplied in
Section~\ref{sec:distribution}. Replace the finite-experiment
justification by that proof, retaining the formal-quotient
qualification. The derivative article's gap was already an
intake warning.

The supporting report, section 4, lines 146 and 148, omits complex
conjugation. Put \(c_q=\gamma+\log(2\pi/q)-1\). The corrected
odd-character identity is
\[
\frac{L'(2,\chi)}{L(2,\chi)}
+\frac12\overline{\frac{L''(-1,\chi)}{L'(-1,\chi)}}=c_q.
\]
For an even primitive character the corrected second logarithmic
identity is
\[
(\log L)''(2,\chi)
-\overline{(\log L)''(-1,\chi)}=1+\frac{\pi^2}{12}.
\]
These are \textbf{established formula errors} for complex
characters, not merely stylistic omissions.
For the odd quartic character modulo five with \(\chi(2)=\ii\),
the uncorrected first formula has a discrepancy approximately
\(-0.1438652751124278\,\ii\).
For the even cubic character modulo seven with
\(\chi(3)=\omega\), the uncorrected second formula has discrepancy
\(-0.0339162597722613\,\ii\).
Section~\ref{sec:bridges} proves the corrected identities at every
order. The corresponding first-level bridges in the two main
Stieltjes articles already have the correct conjugation and should
not be altered on this account.

The claim of mutually independent second-derivative atoms in the
derivative article, lines 303--315, is not established by the PSLQ
experiments. Describe the retained constants as formal generators
unless an arithmetic lower-bound proof is supplied.
The antiderivative article should call its Hurwitz-zeta construction
\emph{a canonical analytic interpolant}, rather than claim uniqueness
from integer data: adding \(c\sin(2\pi z)\) preserves both integer
values and a unit-shift recurrence. Its other small corrections are
listed separately in the register.

\subsection{Polygamma parity and CM evidence}

\noindent\textbf{Sources:}
\path{articles/polylog-polygamma-bridge.tex} and
\path{articles/eisenstein-row-sums-cm-polygamma.tex}.

The Clausen component of the rational grid in the first article,
lines 149--151, has reflection parity \((-1)^{n+1}\) at
polylogarithm weight \(n\). It is not always the odd component.
Thus \(\Cl_2\) is represented by a trigamma difference, while
\(\Cl_3\) uses a tetragamma sum. The displayed laws already have
the correct alternating behavior; the prose needs repair.
The claimed exact converse classification of new Clausen atoms is
not established by the numerical searches. Its positive reductions
and exact \(L\)-function bridges remain useful.

In the CM article, \src{neg:sums} and \src{res:genus} use failed
searches to suggest irrationality or exact algebraic degree.
These are numerical observations. Even an exact quartic equation
only establishes degree at most four until irreducibility and root
identification are supplied.
The weight-two question is resolved in Section~\ref{sec:modular}.
The isolated imaginary trigamma value remains separate from the
modular aggregate identities.

The source's Hurwitz recurrence and its weight-twelve coefficients
are correct in the unnormalized convention:
\[
G_{12}=\frac{18G_4^3+25G_6^2}{143}.
\]
They should not be replaced by normalized Eisenstein coefficients
without carrying the zeta factors.

\subsection{A replacement rule for numerical claims}

Where the only evidence is a finite integer-relation search, use
the statement:
\emph{No relation was found for the specified basis, precision, and
search parameters. The listed constants are retained as formal
generators in this reduction scheme; arithmetic independence is not
established.}

A successful search proposes an identity until an exact proof or
certificate is supplied. A failed search is not automatically a
rigorous bounded-height exclusion either: that needs certified input
enclosures and a certified exclusion argument. A zero canary
coefficient does not prevent a false relation among the other
entries. These distinctions preserve the value of numerical
exploration while preventing unsupported theorem statements.

The older correction
\[
6\Li_2(1/3)-\Li_2(1/9)=\frac{\pi^2}{3}-\log^2 3
\]
to \path{reports/report-1.tex} was already recorded in the
repository README. It is not a new finding of this continuation.
Other pre-existing README caveats remain in force.
The present review does not certify every statement left unchanged.

\clearpage
\phantomsection
\addcontentsline{toc}{section}{References}
\begin{thebibliography}{99}

\bibitem{Repo}
V. Reshetnikov, \emph{ProveIt: Analysis/Polylogarithms/docs},
source snapshot
\path{3a6d80ed618d839915deb0d19ab687f45e81d995}
(October 8, 2026).
\url{https://github.com/VladimirReshetnikov/ProveIt/tree/3a6d80ed618d839915deb0d19ab687f45e81d995/Analysis/Polylogarithms/docs}.
The accompanying source manifest identifies the individual files.

\bibitem{Panzer}
E. Panzer, \emph{The parity theorem for multiple polylogarithms},
Journal of Number Theory \textbf{172} (2017), 93--113.
\url{https://arxiv.org/abs/1512.04482}.
In particular Theorem 1.3, Corollary 1.4, and the depth-two formulas.

\bibitem{XuWang}
C. Xu and W. Wang, \emph{Dirichlet type extensions of Euler sums},
Comptes Rendus Math\'ematique \textbf{361} (2023), 979--1010.
\url{https://doi.org/10.5802/crmath.453}.
The specialization used here is equation (20), section 3.4,
printed page 994.

\bibitem{Kubert}
D. S. Kubert, \emph{The universal ordinary distribution},
Bulletin de la Soci\'et\'e Math\'ematique de France
\textbf{107} (1979), 179--202.
\url{https://doi.org/10.24033/bsmf.1891}.

\bibitem{Ouyang}
Y. Ouyang, \emph{On the universal norm distribution},
Journal of the Ramanujan Mathematical Society
\textbf{17} (2002), no.~4, 287--311.
\url{https://arxiv.org/abs/math/0210297}.

\bibitem{Coffey}
M. W. Coffey, \emph{Series representations for the Stieltjes constants},
Rocky Mountain Journal of Mathematics \textbf{44} (2014), 443--477.
Preprint: \url{https://arxiv.org/abs/0905.1111}.
The parameter-derivative formulas follow from its Hurwitz-zeta
addition and differentiation identities.

\bibitem{KowalskiNikeghbali}
E. Kowalski and A. Nikeghbali,
\emph{Mod-Poisson convergence in probability and number theory},
author manuscript.
\url{https://people.math.ethz.ch/~kowalski/mod-poisson.pdf}.
The permutation-cycle example explains the classical reciprocal-gamma
factor underlying the coefficient mechanism.

\bibitem{RomikScherer}
D. Romik and R. Scherer,
\emph{Alternative summation orders for the Eisenstein series \(G_2\)
and Weierstrass \(\wp\)-function},
arXiv:1811.01523, version 2 (2020).
\url{https://arxiv.org/abs/1811.01523}.

\bibitem{Ramanujan}
S. Ramanujan, \emph{On certain arithmetical functions},
Transactions of the Cambridge Philosophical Society
\textbf{22} (1916), no.~9, 159--184.
Transcription of the original paper:
\url{https://ramanujan.sirinudi.org/Volumes/published/ram18.html}.

\bibitem{ABP}
G. W. Anderson, W. D. Brownawell, and M. A. Papanikolas,
\emph{Determination of the algebraic relations among special
\(\Gamma\)-values in positive characteristic},
Annals of Mathematics \textbf{160} (2004), 237--313.
\url{https://annals.math.princeton.edu/wp-content/uploads/annals-v160-n1-p06.pdf}.
Section 1.4.2 distinguishes the classical Koblitz--Ogus sufficient
criterion from Rohrlich's conjectural converse; the paper's
positive-characteristic theorem is a different result.

\bibitem{DLMF}
NIST, \emph{Digital Library of Mathematical Functions},
sections 5.11 (gamma asymptotics), 23 (Weierstrass elliptic
functions), 25.11 (Hurwitz zeta), and 25.15 (Dirichlet \(L\)-functions).
\url{https://dlmf.nist.gov/}.
Consulted October 8, 2026.

\end{thebibliography}

\end{document}
